%%%%%%%%%%%%%%%%%%%%%%%%%%%%%%%%%%%%%%%%%%%%%%%%%%%%%%%%%%%%
% 9.1 ORDINARY DIFFERENTIAL EQUATIONS
% The Composition of Advanced Mathematics
%%%%%%%%%%%%%%%%%%%%%%%%%%%%%%%%%%%%%%%%%%%%%%%%%%%%%%%%%%%%

\section{Ordinary Differential Equations}
\label{sec:ordinary-differential-equations}

\prerequisite{
Functions, limits, derivatives, integrals, sequences and series,
linear algebra, matrices, eigenvalues and eigenvectors, complex numbers,
and basic mathematical modeling.
}

\learningobjectives{
By the end of this section, the reader should be able to:
\begin{itemize}
    \item define an ordinary differential equation;
    \item distinguish ordinary and partial differential equations;
    \item determine the order and degree of a differential equation;
    \item distinguish linear and nonlinear differential equations;
    \item distinguish autonomous and nonautonomous equations;
    \item distinguish general, particular, and singular solutions;
    \item understand initial-value and boundary-value problems;
    \item verify proposed solutions;
    \item understand existence and uniqueness conceptually;
    \item interpret direction fields;
    \item analyze equilibrium solutions and stability in one dimension;
    \item solve separable first-order equations;
    \item solve first-order linear equations using integrating factors;
    \item solve exact differential equations;
    \item understand integrating factors for nonexact equations
          conceptually;
    \item solve Bernoulli equations;
    \item solve common homogeneous first-order equations;
    \item model exponential growth and decay;
    \item model logistic growth;
    \item model Newton's law of cooling;
    \item model mixing problems;
    \item solve second-order linear homogeneous equations with constant
          coefficients;
    \item handle distinct, repeated, and complex characteristic roots;
    \item solve nonhomogeneous linear equations;
    \item use undetermined coefficients;
    \item use variation of parameters conceptually and computationally;
    \item understand the superposition principle;
    \item compute and interpret the Wronskian;
    \item reduce higher-order equations to first-order systems;
    \item solve linear systems using eigenvalues and eigenvectors;
    \item classify planar linear systems;
    \item understand fundamental matrices;
    \item analyze mechanical vibration equations;
    \item distinguish free, damped, and forced oscillation;
    \item understand resonance;
    \item recognize nonlinear first-order and second-order systems;
    \item understand numerical approximation by Euler's method;
    \item recognize improved Euler and Runge--Kutta methods;
    \item connect ordinary differential equations with dynamical systems,
          modeling, numerical analysis, and transform methods.
\end{itemize}
}

%%%%%%%%%%%%%%%%%%%%%%%%%%%%%%%%%%%%%%%%%%%%%%%%%%%%%%%%%%%%
\subsection{What Is a Differential Equation?}
\label{subsec:what-is-differential-equation}
%%%%%%%%%%%%%%%%%%%%%%%%%%%%%%%%%%%%%%%%%%%%%%%%%%%%%%%%%%%%

A differential equation is an equation involving an unknown function and
one or more of its derivatives.

For example,

\[
\boxed{
\frac{dy}{dx}
=
3x^2
}
\]

relates the unknown function \(y(x)\) to its derivative.

Another example is

\[
\boxed{
y''+4y=0.
}
\]

The goal is to determine functions satisfying the equation.

%%%%%%%%%%%%%%%%%%%%%%%%%%%%%%%%%%%%%%%%%%%%%%%%%%%%%%%%%%%%
\subsection{Ordinary Differential Equations}
\label{subsec:definition-ode}
%%%%%%%%%%%%%%%%%%%%%%%%%%%%%%%%%%%%%%%%%%%%%%%%%%%%%%%%%%%%

\begin{definitionbox}{Ordinary Differential Equation}
An ordinary differential equation, abbreviated ODE, is a differential
equation involving an unknown function of one independent variable and
its ordinary derivatives.
\end{definitionbox}

A general ODE may be written

\[
\boxed{
F
\left(
x,
y,
y',
y'',
\ldots,
y^{(n)}
\right)
=
0.
}
\]

%%%%%%%%%%%%%%%%%%%%%%%%%%%%%%%%%%%%%%%%%%%%%%%%%%%%%%%%%%%%
\subsection{ODE Versus PDE}
\label{subsec:ode-versus-pde}
%%%%%%%%%%%%%%%%%%%%%%%%%%%%%%%%%%%%%%%%%%%%%%%%%%%%%%%%%%%%

An ODE involves derivatives with respect to one independent variable.

For example,

\[
\boxed{
y''+y=0.
}
\]

A partial differential equation involves partial derivatives with respect
to multiple independent variables.

For example,

\[
\boxed{
\frac{\partial u}{\partial t}
=
k
\frac{\partial^2u}{\partial x^2}.
}
\]

Partial differential equations are developed later in Chapter 9.

%%%%%%%%%%%%%%%%%%%%%%%%%%%%%%%%%%%%%%%%%%%%%%%%%%%%%%%%%%%%
\subsection{Order of a Differential Equation}
\label{subsec:ode-order}
%%%%%%%%%%%%%%%%%%%%%%%%%%%%%%%%%%%%%%%%%%%%%%%%%%%%%%%%%%%%

\begin{definitionbox}{Order}
The order of a differential equation is the highest derivative appearing
in the equation.
\end{definitionbox}

For example,

\[
y'''+2y'-y=0
\]

is a third-order differential equation.

Thus,

\[
\boxed{
\operatorname{order}=3.
}
\]

%%%%%%%%%%%%%%%%%%%%%%%%%%%%%%%%%%%%%%%%%%%%%%%%%%%%%%%%%%%%
\subsection{Degree of a Differential Equation}
\label{subsec:ode-degree}
%%%%%%%%%%%%%%%%%%%%%%%%%%%%%%%%%%%%%%%%%%%%%%%%%%%%%%%%%%%%

When a differential equation is polynomial in its derivatives, its
degree is the power of the highest-order derivative after radicals and
fractions involving derivatives have been removed.

For example,

\[
\boxed{
(y'')^2+y'=x
}
\]

has order \(2\) and degree \(2\).

Degree is not defined in this elementary sense for every differential
equation.

%%%%%%%%%%%%%%%%%%%%%%%%%%%%%%%%%%%%%%%%%%%%%%%%%%%%%%%%%%%%
\subsection{Linear Differential Equations}
\label{subsec:linear-odes}
%%%%%%%%%%%%%%%%%%%%%%%%%%%%%%%%%%%%%%%%%%%%%%%%%%%%%%%%%%%%

An \(n\)-th order linear ODE has the form

\[
\boxed{
a_n(x)y^{(n)}
+
a_{n-1}(x)y^{(n-1)}
+
\cdots
+
a_1(x)y'
+
a_0(x)y
=
g(x).
}
\]

The unknown function and its derivatives appear only to the first power
and are not multiplied together.

%%%%%%%%%%%%%%%%%%%%%%%%%%%%%%%%%%%%%%%%%%%%%%%%%%%%%%%%%%%%
\subsection{Homogeneous Linear Equations}
\label{subsec:homogeneous-linear-odes}
%%%%%%%%%%%%%%%%%%%%%%%%%%%%%%%%%%%%%%%%%%%%%%%%%%%%%%%%%%%%

A linear equation is homogeneous when

\[
\boxed{
g(x)=0.
}
\]

Thus,

\[
\boxed{
a_n(x)y^{(n)}
+
\cdots
+
a_0(x)y
=
0.
}
\]

For example,

\[
\boxed{
y''-3y'+2y=0.
}
\]

%%%%%%%%%%%%%%%%%%%%%%%%%%%%%%%%%%%%%%%%%%%%%%%%%%%%%%%%%%%%
\subsection{Nonhomogeneous Linear Equations}
\label{subsec:nonhomogeneous-linear-odes}
%%%%%%%%%%%%%%%%%%%%%%%%%%%%%%%%%%%%%%%%%%%%%%%%%%%%%%%%%%%%

A linear equation is nonhomogeneous when

\[
\boxed{
g(x)\neq0.
}
\]

For example,

\[
\boxed{
y''+4y=\sin x.
}
\]

The function \(g(x)\) is often called the forcing term.

%%%%%%%%%%%%%%%%%%%%%%%%%%%%%%%%%%%%%%%%%%%%%%%%%%%%%%%%%%%%
\subsection{Nonlinear Differential Equations}
\label{subsec:nonlinear-odes}
%%%%%%%%%%%%%%%%%%%%%%%%%%%%%%%%%%%%%%%%%%%%%%%%%%%%%%%%%%%%

An ODE is nonlinear if the unknown function or its derivatives appear
nonlinearly.

Examples include

\[
\boxed{
y'=y^2-x,
}
\]

\[
\boxed{
yy''+(y')^2=0,
}
\]

and

\[
\boxed{
y''+\sin y=0.
}
\]

Nonlinear equations can behave very differently from linear equations.

%%%%%%%%%%%%%%%%%%%%%%%%%%%%%%%%%%%%%%%%%%%%%%%%%%%%%%%%%%%%
\subsection{Autonomous Equations}
\label{subsec:autonomous-odes}
%%%%%%%%%%%%%%%%%%%%%%%%%%%%%%%%%%%%%%%%%%%%%%%%%%%%%%%%%%%%

A first-order autonomous equation has the form

\[
\boxed{
y'
=
f(y).
}
\]

The independent variable does not appear explicitly.

For example,

\[
\boxed{
y'
=
y(1-y).
}
\]

Autonomous equations are especially important in dynamical systems.

%%%%%%%%%%%%%%%%%%%%%%%%%%%%%%%%%%%%%%%%%%%%%%%%%%%%%%%%%%%%
\subsection{Nonautonomous Equations}
\label{subsec:nonautonomous-odes}
%%%%%%%%%%%%%%%%%%%%%%%%%%%%%%%%%%%%%%%%%%%%%%%%%%%%%%%%%%%%

A nonautonomous equation depends explicitly on the independent variable.

For example,

\[
\boxed{
y'
=
x+y.
}
\]

Here the evolution depends directly on \(x\) as well as on the state
\(y\).

%%%%%%%%%%%%%%%%%%%%%%%%%%%%%%%%%%%%%%%%%%%%%%%%%%%%%%%%%%%%
\subsection{Solutions of Differential Equations}
\label{subsec:ode-solutions}
%%%%%%%%%%%%%%%%%%%%%%%%%%%%%%%%%%%%%%%%%%%%%%%%%%%%%%%%%%%%

A function \(y(x)\) is a solution if substituting it and its derivatives
into the differential equation makes the equation true throughout the
relevant interval.

For

\[
y'=2x,
\]

a family of solutions is

\[
\boxed{
y=x^2+C.
}
\]

%%%%%%%%%%%%%%%%%%%%%%%%%%%%%%%%%%%%%%%%%%%%%%%%%%%%%%%%%%%%
\subsection{General Solution}
\label{subsec:general-solution-ode}
%%%%%%%%%%%%%%%%%%%%%%%%%%%%%%%%%%%%%%%%%%%%%%%%%%%%%%%%%%%%

A general solution usually contains arbitrary constants.

For a first-order equation, one commonly expects one constant.

For an \(n\)-th order equation, one often expects \(n\) independent
constants.

For example,

\[
\boxed{
y''+y=0
}
\]

has general solution

\[
\boxed{
y
=
C_1\cos x
+
C_2\sin x.
}
\]

%%%%%%%%%%%%%%%%%%%%%%%%%%%%%%%%%%%%%%%%%%%%%%%%%%%%%%%%%%%%
\subsection{Particular Solution}
\label{subsec:particular-solution-ode}
%%%%%%%%%%%%%%%%%%%%%%%%%%%%%%%%%%%%%%%%%%%%%%%%%%%%%%%%%%%%

A particular solution is obtained after choosing specific values of the
arbitrary constants.

For example, from

\[
y=x^2+C,
\]

choosing

\[
C=3
\]

gives

\[
\boxed{
y=x^2+3.
}
\]

%%%%%%%%%%%%%%%%%%%%%%%%%%%%%%%%%%%%%%%%%%%%%%%%%%%%%%%%%%%%
\subsection{Singular Solutions}
\label{subsec:singular-solutions-ode}
%%%%%%%%%%%%%%%%%%%%%%%%%%%%%%%%%%%%%%%%%%%%%%%%%%%%%%%%%%%%

Some nonlinear differential equations possess solutions that cannot be
obtained by choosing constants in the general solution family.

These are called singular solutions.

They require separate analysis and often arise when algebraic operations
used during solution procedures remove special cases.

%%%%%%%%%%%%%%%%%%%%%%%%%%%%%%%%%%%%%%%%%%%%%%%%%%%%%%%%%%%%
\subsection{Initial-Value Problems}
\label{subsec:initial-value-problems}
%%%%%%%%%%%%%%%%%%%%%%%%%%%%%%%%%%%%%%%%%%%%%%%%%%%%%%%%%%%%

\begin{definitionbox}{Initial-Value Problem}
An initial-value problem, abbreviated IVP, consists of a differential
equation together with conditions specified at a single value of the
independent variable.
\end{definitionbox}

For a first-order equation,

\[
\boxed{
y'=f(x,y),
\qquad
y(x_0)=y_0.
}
\]

For a second-order equation,

\[
\boxed{
y''=f(x,y,y'),
}
\]

with

\[
\boxed{
y(x_0)=y_0,
\qquad
y'(x_0)=v_0.
}
\]

%%%%%%%%%%%%%%%%%%%%%%%%%%%%%%%%%%%%%%%%%%%%%%%%%%%%%%%%%%%%
\subsection{Boundary-Value Problems}
\label{subsec:boundary-value-problems}
%%%%%%%%%%%%%%%%%%%%%%%%%%%%%%%%%%%%%%%%%%%%%%%%%%%%%%%%%%%%

A boundary-value problem specifies conditions at different values of the
independent variable.

For example,

\[
\boxed{
y''+\lambda y=0,
}
\]

with

\[
\boxed{
y(0)=0,
\qquad
y(L)=0.
}
\]

Boundary-value problems arise naturally in mechanics, heat flow, and
eigenvalue problems.

%%%%%%%%%%%%%%%%%%%%%%%%%%%%%%%%%%%%%%%%%%%%%%%%%%%%%%%%%%%%
\subsection{Verifying a Solution}
\label{subsec:verify-ode-solution}
%%%%%%%%%%%%%%%%%%%%%%%%%%%%%%%%%%%%%%%%%%%%%%%%%%%%%%%%%%%%

To verify a proposed solution:

\begin{enumerate}
    \item compute the required derivatives;
    \item substitute into the differential equation;
    \item simplify;
    \item check all initial or boundary conditions.
\end{enumerate}

%%%%%%%%%%%%%%%%%%%%%%%%%%%%%%%%%%%%%%%%%%%%%%%%%%%%%%%%%%%%
\subsection{Worked Example: Verifying a Solution}
\label{subsec:worked-verify-ode}
%%%%%%%%%%%%%%%%%%%%%%%%%%%%%%%%%%%%%%%%%%%%%%%%%%%%%%%%%%%%

\begin{workedexamplebox}{Verification}

Verify that

\[
y=e^{2x}
\]

solves

\[
y'-2y=0.
\]

Differentiate:

\[
y'
=
2e^{2x}.
\]

Substitute:

\[
y'-2y
=
2e^{2x}
-
2e^{2x}
=
0.
\]

\begin{answerbox}
\[
\boxed{
y=e^{2x}
\text{ is a solution.}
}
\]
\end{answerbox}

\end{workedexamplebox}

%%%%%%%%%%%%%%%%%%%%%%%%%%%%%%%%%%%%%%%%%%%%%%%%%%%%%%%%%%%%
\subsection{Existence and Uniqueness}
\label{subsec:ode-existence-uniqueness}
%%%%%%%%%%%%%%%%%%%%%%%%%%%%%%%%%%%%%%%%%%%%%%%%%%%%%%%%%%%%

An important theoretical question is whether an IVP has:

\[
\boxed{
\text{a solution},
}
\]

and whether that solution is

\[
\boxed{
\text{unique}.
}
\]

These are separate questions.

An equation may have no solution, one solution, or multiple solutions
through a given initial condition.

%%%%%%%%%%%%%%%%%%%%%%%%%%%%%%%%%%%%%%%%%%%%%%%%%%%%%%%%%%%%
\subsection{Picard--Lindelöf Theorem}
\label{subsec:picard-lindelof}
%%%%%%%%%%%%%%%%%%%%%%%%%%%%%%%%%%%%%%%%%%%%%%%%%%%%%%%%%%%%

\begin{theorem}[Local Existence and Uniqueness]
Consider

\[
y'=f(x,y),
\qquad
y(x_0)=y_0.
\]

If \(f\) is continuous near \((x_0,y_0)\) and is locally Lipschitz in
\(y\), then there exists an interval around \(x_0\) on which the IVP has
a unique solution.
\end{theorem}

A common sufficient condition is continuity of

\[
\frac{\partial f}{\partial y}
\]

near the initial point.

%%%%%%%%%%%%%%%%%%%%%%%%%%%%%%%%%%%%%%%%%%%%%%%%%%%%%%%%%%%%
\subsection{Failure of Uniqueness}
\label{subsec:failure-uniqueness-ode}
%%%%%%%%%%%%%%%%%%%%%%%%%%%%%%%%%%%%%%%%%%%%%%%%%%%%%%%%%%%%

Consider

\[
\boxed{
y'
=
3y^{2/3},
\qquad
y(0)=0.
}
\]

The function

\[
f(y)=3y^{2/3}
\]

is continuous, but its derivative with respect to \(y\) is problematic at
\(y=0\).

The IVP can possess multiple solutions.

Thus continuity alone does not guarantee uniqueness.

%%%%%%%%%%%%%%%%%%%%%%%%%%%%%%%%%%%%%%%%%%%%%%%%%%%%%%%%%%%%
\subsection{Direction Fields}
\label{subsec:direction-fields}
%%%%%%%%%%%%%%%%%%%%%%%%%%%%%%%%%%%%%%%%%%%%%%%%%%%%%%%%%%%%

For

\[
y'=f(x,y),
\]

a direction field assigns slope

\[
\boxed{
f(x,y)
}
\]

to each point \((x,y)\).

Small line segments visualize how solution curves move through the plane.

Direction fields are useful when an exact formula is unavailable.

%%%%%%%%%%%%%%%%%%%%%%%%%%%%%%%%%%%%%%%%%%%%%%%%%%%%%%%%%%%%
\subsection{Equilibrium Solutions}
\label{subsec:equilibrium-solutions}
%%%%%%%%%%%%%%%%%%%%%%%%%%%%%%%%%%%%%%%%%%%%%%%%%%%%%%%%%%%%

For an autonomous equation

\[
y'=f(y),
\]

an equilibrium is a constant value

\[
\boxed{
y=y^*
}
\]

satisfying

\[
\boxed{
f(y^*)=0.
}
\]

Then

\[
y'=0,
\]

so the solution remains constant.

%%%%%%%%%%%%%%%%%%%%%%%%%%%%%%%%%%%%%%%%%%%%%%%%%%%%%%%%%%%%
\subsection{Phase Line}
\label{subsec:phase-line}
%%%%%%%%%%%%%%%%%%%%%%%%%%%%%%%%%%%%%%%%%%%%%%%%%%%%%%%%%%%%

A one-dimensional autonomous equation

\[
y'=f(y)
\]

can be analyzed using a phase line.

Intervals are marked according to the sign of \(f(y)\):

\[
\boxed{
f(y)>0
\Rightarrow
y\text{ increases},
}
\]

\[
\boxed{
f(y)<0
\Rightarrow
y\text{ decreases}.
}
\]

%%%%%%%%%%%%%%%%%%%%%%%%%%%%%%%%%%%%%%%%%%%%%%%%%%%%%%%%%%%%
\subsection{Stable Equilibrium}
\label{subsec:stable-equilibrium-one-dimensional}
%%%%%%%%%%%%%%%%%%%%%%%%%%%%%%%%%%%%%%%%%%%%%%%%%%%%%%%%%%%%

An equilibrium is locally stable if nearby solutions remain near it.

It is asymptotically stable if nearby solutions approach it as time
increases.

In one dimension, arrows pointing toward an equilibrium from both sides
indicate asymptotic stability.

%%%%%%%%%%%%%%%%%%%%%%%%%%%%%%%%%%%%%%%%%%%%%%%%%%%%%%%%%%%%
\subsection{Unstable Equilibrium}
\label{subsec:unstable-equilibrium-one-dimensional}
%%%%%%%%%%%%%%%%%%%%%%%%%%%%%%%%%%%%%%%%%%%%%%%%%%%%%%%%%%%%

An equilibrium is unstable if arbitrarily nearby solutions move away from
it.

On a phase line, arrows pointing away from the equilibrium on both sides
indicate instability.

%%%%%%%%%%%%%%%%%%%%%%%%%%%%%%%%%%%%%%%%%%%%%%%%%%%%%%%%%%%%
\subsection{Derivative Stability Test}
\label{subsec:one-dimensional-stability-test}
%%%%%%%%%%%%%%%%%%%%%%%%%%%%%%%%%%%%%%%%%%%%%%%%%%%%%%%%%%%%

For

\[
y'=f(y),
\]

suppose

\[
f(y^*)=0.
\]

If

\[
\boxed{
f'(y^*)<0,
}
\]

the equilibrium is locally asymptotically stable.

If

\[
\boxed{
f'(y^*)>0,
}
\]

the equilibrium is unstable.

If

\[
f'(y^*)=0,
\]

the linear test is inconclusive.

%%%%%%%%%%%%%%%%%%%%%%%%%%%%%%%%%%%%%%%%%%%%%%%%%%%%%%%%%%%%
\subsection{Separable Differential Equations}
\label{subsec:separable-equations}
%%%%%%%%%%%%%%%%%%%%%%%%%%%%%%%%%%%%%%%%%%%%%%%%%%%%%%%%%%%%

A first-order equation is separable if it can be written

\[
\boxed{
\frac{dy}{dx}
=
g(x)h(y).
}
\]

Rearrange:

\[
\boxed{
\frac1{h(y)}
\,dy
=
g(x)\,dx.
}
\]

Then integrate both sides:

\[
\boxed{
\int
\frac1{h(y)}
\,dy
=
\int
g(x)\,dx.
}
\]

%%%%%%%%%%%%%%%%%%%%%%%%%%%%%%%%%%%%%%%%%%%%%%%%%%%%%%%%%%%%
\subsection{Separation Caveat}
\label{subsec:separation-caveat}
%%%%%%%%%%%%%%%%%%%%%%%%%%%%%%%%%%%%%%%%%%%%%%%%%%%%%%%%%%%%

When dividing by

\[
h(y),
\]

one must first check whether

\[
\boxed{
h(y)=0
}
\]

produces equilibrium solutions.

Otherwise those solutions can be lost during division.

%%%%%%%%%%%%%%%%%%%%%%%%%%%%%%%%%%%%%%%%%%%%%%%%%%%%%%%%%%%%
\subsection{Worked Example: Separable Equation}
\label{subsec:worked-separable-ode}
%%%%%%%%%%%%%%%%%%%%%%%%%%%%%%%%%%%%%%%%%%%%%%%%%%%%%%%%%%%%

\begin{workedexamplebox}{Solving a Separable Equation}

Solve

\[
\frac{dy}{dx}
=
xy.
\]

Separate variables:

\[
\frac1y\,dy
=
x\,dx.
\]

Integrate:

\[
\ln|y|
=
\frac{x^2}{2}
+
C.
\]

Exponentiate:

\[
|y|
=
e^C
e^{x^2/2}.
\]

Absorb the sign and constant into \(C_1\):

\begin{answerbox}
\[
\boxed{
y
=
Ce^{x^2/2}.
}
\]
\end{answerbox}

The solution \(y=0\) is included by \(C=0\).

\end{workedexamplebox}

%%%%%%%%%%%%%%%%%%%%%%%%%%%%%%%%%%%%%%%%%%%%%%%%%%%%%%%%%%%%
\subsection{Initial Condition for a Separable Equation}
\label{subsec:separable-ivp}
%%%%%%%%%%%%%%%%%%%%%%%%%%%%%%%%%%%%%%%%%%%%%%%%%%%%%%%%%%%%

Suppose

\[
y'
=
xy,
\qquad
y(0)=3.
\]

From

\[
y
=
Ce^{x^2/2},
\]

apply the initial condition:

\[
3
=
Ce^0.
\]

Thus

\[
\boxed{
C=3.
}
\]

Therefore,

\[
\boxed{
y
=
3e^{x^2/2}.
}
\]

%%%%%%%%%%%%%%%%%%%%%%%%%%%%%%%%%%%%%%%%%%%%%%%%%%%%%%%%%%%%
\subsection{First-Order Linear Equations}
\label{subsec:first-order-linear-odes}
%%%%%%%%%%%%%%%%%%%%%%%%%%%%%%%%%%%%%%%%%%%%%%%%%%%%%%%%%%%%

A first-order linear ODE has standard form

\[
\boxed{
y'
+
P(x)y
=
Q(x).
}
\]

The integrating-factor method converts the left side into the derivative
of a product.

%%%%%%%%%%%%%%%%%%%%%%%%%%%%%%%%%%%%%%%%%%%%%%%%%%%%%%%%%%%%
\subsection{Integrating Factor}
\label{subsec:integrating-factor}
%%%%%%%%%%%%%%%%%%%%%%%%%%%%%%%%%%%%%%%%%%%%%%%%%%%%%%%%%%%%

Define

\[
\boxed{
\mu(x)
=
e^{\int P(x)\,dx}.
}
\]

Multiply the differential equation by \(\mu(x)\):

\[
\mu y'
+
\mu P y
=
\mu Q.
\]

Because

\[
\mu'
=
P\mu,
\]

the left side becomes

\[
\boxed{
(\mu y)'.
}
\]

Thus,

\[
\boxed{
(\mu y)'
=
\mu Q.
}
\]

%%%%%%%%%%%%%%%%%%%%%%%%%%%%%%%%%%%%%%%%%%%%%%%%%%%%%%%%%%%%
\subsection{General First-Order Linear Solution}
\label{subsec:first-order-linear-general-solution}
%%%%%%%%%%%%%%%%%%%%%%%%%%%%%%%%%%%%%%%%%%%%%%%%%%%%%%%%%%%%

Integrating,

\[
\mu(x)y
=
\int
\mu(x)Q(x)\,dx
+
C.
\]

Therefore,

\[
\boxed{
y
=
\frac1{\mu(x)}
\left[
\int
\mu(x)Q(x)\,dx
+
C
\right].
}
\]

%%%%%%%%%%%%%%%%%%%%%%%%%%%%%%%%%%%%%%%%%%%%%%%%%%%%%%%%%%%%
\subsection{Worked Example: Linear First-Order Equation}
\label{subsec:worked-linear-first-order}
%%%%%%%%%%%%%%%%%%%%%%%%%%%%%%%%%%%%%%%%%%%%%%%%%%%%%%%%%%%%

\begin{workedexamplebox}{Integrating Factor}

Solve

\[
y'+2y=e^x.
\]

Here,

\[
P(x)=2.
\]

The integrating factor is

\[
\mu(x)
=
e^{\int2\,dx}
=
e^{2x}.
\]

Multiply through:

\[
e^{2x}y'
+
2e^{2x}y
=
e^{3x}.
\]

Therefore,

\[
(e^{2x}y)'
=
e^{3x}.
\]

Integrate:

\[
e^{2x}y
=
\frac13e^{3x}
+
C.
\]

Thus,

\begin{answerbox}
\[
\boxed{
y
=
\frac13e^x
+
Ce^{-2x}.
}
\]
\end{answerbox}

\end{workedexamplebox}

%%%%%%%%%%%%%%%%%%%%%%%%%%%%%%%%%%%%%%%%%%%%%%%%%%%%%%%%%%%%
\subsection{Exact Differential Equations}
\label{subsec:exact-differential-equations}
%%%%%%%%%%%%%%%%%%%%%%%%%%%%%%%%%%%%%%%%%%%%%%%%%%%%%%%%%%%%

Consider an equation

\[
\boxed{
M(x,y)\,dx
+
N(x,y)\,dy
=
0.
}
\]

It is exact if there exists a potential function

\[
F(x,y)
\]

such that

\[
\boxed{
F_x=M,
\qquad
F_y=N.
}
\]

Then the solution is

\[
\boxed{
F(x,y)=C.
}
\]

%%%%%%%%%%%%%%%%%%%%%%%%%%%%%%%%%%%%%%%%%%%%%%%%%%%%%%%%%%%%
\subsection{Exactness Test}
\label{subsec:exactness-test}
%%%%%%%%%%%%%%%%%%%%%%%%%%%%%%%%%%%%%%%%%%%%%%%%%%%%%%%%%%%%

If the relevant partial derivatives are continuous, exactness requires

\[
\boxed{
\frac{\partial M}{\partial y}
=
\frac{\partial N}{\partial x}.
}
\]

This follows from equality of mixed partial derivatives:

\[
F_{xy}=F_{yx}.
\]

%%%%%%%%%%%%%%%%%%%%%%%%%%%%%%%%%%%%%%%%%%%%%%%%%%%%%%%%%%%%
\subsection{Worked Example: Exact Equation}
\label{subsec:worked-exact-equation}
%%%%%%%%%%%%%%%%%%%%%%%%%%%%%%%%%%%%%%%%%%%%%%%%%%%%%%%%%%%%

\begin{workedexamplebox}{Potential Function}

Solve

\[
(2x+y)\,dx
+
(x+2y)\,dy
=
0.
\]

Here,

\[
M=2x+y,
\qquad
N=x+2y.
\]

Check:

\[
M_y=1,
\qquad
N_x=1.
\]

So the equation is exact.

Integrate \(M\) with respect to \(x\):

\[
F
=
x^2+xy+g(y).
\]

Differentiate with respect to \(y\):

\[
F_y
=
x+g'(y).
\]

Set equal to \(N\):

\[
x+g'(y)
=
x+2y.
\]

Hence

\[
g'(y)=2y,
\]

so

\[
g(y)=y^2.
\]

Therefore,

\begin{answerbox}
\[
\boxed{
x^2+xy+y^2=C.
}
\]
\end{answerbox}

\end{workedexamplebox}

%%%%%%%%%%%%%%%%%%%%%%%%%%%%%%%%%%%%%%%%%%%%%%%%%%%%%%%%%%%%
\subsection{Integrating Factors for Nonexact Equations}
\label{subsec:integrating-factors-nonexact}
%%%%%%%%%%%%%%%%%%%%%%%%%%%%%%%%%%%%%%%%%%%%%%%%%%%%%%%%%%%%

A nonexact equation

\[
M\,dx+N\,dy=0
\]

may sometimes become exact after multiplication by an integrating factor

\[
\boxed{
\mu(x,y).
}
\]

That is,

\[
\boxed{
\mu M\,dx
+
\mu N\,dy
=
0
}
\]

is exact.

Finding such an integrating factor can be difficult in general.

%%%%%%%%%%%%%%%%%%%%%%%%%%%%%%%%%%%%%%%%%%%%%%%%%%%%%%%%%%%%
\subsection{Bernoulli Equations}
\label{subsec:bernoulli-equations}
%%%%%%%%%%%%%%%%%%%%%%%%%%%%%%%%%%%%%%%%%%%%%%%%%%%%%%%%%%%%

A Bernoulli equation has form

\[
\boxed{
y'
+
P(x)y
=
Q(x)y^n,
}
\]

where

\[
n\neq0,1.
\]

The substitution

\[
\boxed{
v
=
y^{1-n}
}
\]

transforms the equation into a first-order linear ODE in \(v\).

%%%%%%%%%%%%%%%%%%%%%%%%%%%%%%%%%%%%%%%%%%%%%%%%%%%%%%%%%%%%
\subsection{Deriving the Bernoulli Transformation}
\label{subsec:bernoulli-transformation}
%%%%%%%%%%%%%%%%%%%%%%%%%%%%%%%%%%%%%%%%%%%%%%%%%%%%%%%%%%%%

Divide by \(y^n\):

\[
y^{-n}y'
+
P(x)y^{1-n}
=
Q(x).
\]

Let

\[
v=y^{1-n}.
\]

Then

\[
v'
=
(1-n)y^{-n}y'.
\]

Therefore,

\[
\boxed{
v'
+
(1-n)P(x)v
=
(1-n)Q(x).
}
\]

This is linear in \(v\).

%%%%%%%%%%%%%%%%%%%%%%%%%%%%%%%%%%%%%%%%%%%%%%%%%%%%%%%%%%%%
\subsection{Homogeneous First-Order Substitution}
\label{subsec:homogeneous-first-order-substitution}
%%%%%%%%%%%%%%%%%%%%%%%%%%%%%%%%%%%%%%%%%%%%%%%%%%%%%%%%%%%%

Some first-order equations have form

\[
\boxed{
y'
=
F
\left(
\frac{y}{x}
\right).
}
\]

Use

\[
\boxed{
v
=
\frac{y}{x},
\qquad
y=vx.
}
\]

Then

\[
\boxed{
y'
=
v
+
xv'.
}
\]

The transformed equation is often separable.

%%%%%%%%%%%%%%%%%%%%%%%%%%%%%%%%%%%%%%%%%%%%%%%%%%%%%%%%%%%%
\subsection{Exponential Growth and Decay}
\label{subsec:exponential-growth-decay}
%%%%%%%%%%%%%%%%%%%%%%%%%%%%%%%%%%%%%%%%%%%%%%%%%%%%%%%%%%%%

A basic model assumes the rate of change is proportional to the current
amount:

\[
\boxed{
\frac{dP}{dt}
=
kP.
}
\]

Separate:

\[
\frac{dP}{P}
=
k\,dt.
\]

Integrating gives

\[
\boxed{
P(t)
=
P_0e^{kt}.
}
\]

%%%%%%%%%%%%%%%%%%%%%%%%%%%%%%%%%%%%%%%%%%%%%%%%%%%%%%%%%%%%
\subsection{Growth Versus Decay}
\label{subsec:growth-versus-decay}
%%%%%%%%%%%%%%%%%%%%%%%%%%%%%%%%%%%%%%%%%%%%%%%%%%%%%%%%%%%%

If

\[
\boxed{
k>0,
}
\]

the model describes exponential growth.

If

\[
\boxed{
k<0,
}
\]

it describes exponential decay.

%%%%%%%%%%%%%%%%%%%%%%%%%%%%%%%%%%%%%%%%%%%%%%%%%%%%%%%%%%%%
\subsection{Doubling Time}
\label{subsec:doubling-time}
%%%%%%%%%%%%%%%%%%%%%%%%%%%%%%%%%%%%%%%%%%%%%%%%%%%%%%%%%%%%

For exponential growth,

\[
P(t)=P_0e^{kt}.
\]

The doubling time \(T_d\) satisfies

\[
2P_0
=
P_0e^{kT_d}.
\]

Therefore,

\[
\boxed{
T_d
=
\frac{
\ln2
}{
k
}.
}
\]

%%%%%%%%%%%%%%%%%%%%%%%%%%%%%%%%%%%%%%%%%%%%%%%%%%%%%%%%%%%%
\subsection{Half-Life}
\label{subsec:half-life}
%%%%%%%%%%%%%%%%%%%%%%%%%%%%%%%%%%%%%%%%%%%%%%%%%%%%%%%%%%%%

For exponential decay,

\[
P(t)=P_0e^{-kt},
\qquad
k>0.
\]

The half-life satisfies

\[
\frac12P_0
=
P_0e^{-kT_{1/2}}.
\]

Thus,

\[
\boxed{
T_{1/2}
=
\frac{
\ln2
}{
k
}.
}
\]

%%%%%%%%%%%%%%%%%%%%%%%%%%%%%%%%%%%%%%%%%%%%%%%%%%%%%%%%%%%%
\subsection{Logistic Growth}
\label{subsec:logistic-growth-ode}
%%%%%%%%%%%%%%%%%%%%%%%%%%%%%%%%%%%%%%%%%%%%%%%%%%%%%%%%%%%%

Exponential growth is unrealistic when resources are limited.

The logistic equation is

\[
\boxed{
\frac{dP}{dt}
=
rP
\left(
1-\frac{P}{K}
\right),
}
\]

where:

\[
r>0
=
\text{intrinsic growth rate},
\]

and

\[
K>0
=
\text{carrying capacity}.
\]

%%%%%%%%%%%%%%%%%%%%%%%%%%%%%%%%%%%%%%%%%%%%%%%%%%%%%%%%%%%%
\subsection{Logistic Equilibria}
\label{subsec:logistic-equilibria}
%%%%%%%%%%%%%%%%%%%%%%%%%%%%%%%%%%%%%%%%%%%%%%%%%%%%%%%%%%%%

Set

\[
P'=0.
\]

Then

\[
rP
\left(
1-\frac{P}{K}
\right)
=
0.
\]

Thus the equilibria are

\[
\boxed{
P=0
}
\]

and

\[
\boxed{
P=K.
}
\]

For positive populations,

\[
P=K
\]

is asymptotically stable.

%%%%%%%%%%%%%%%%%%%%%%%%%%%%%%%%%%%%%%%%%%%%%%%%%%%%%%%%%%%%
\subsection{Logistic Solution}
\label{subsec:logistic-solution}
%%%%%%%%%%%%%%%%%%%%%%%%%%%%%%%%%%%%%%%%%%%%%%%%%%%%%%%%%%%%

The logistic equation has solution

\[
\boxed{
P(t)
=
\frac{
K
}{
1+Ae^{-rt}
},
}
\]

where \(A\) is determined by the initial condition.

If

\[
P(0)=P_0,
\]

then

\[
\boxed{
A
=
\frac{
K-P_0
}{
P_0
}.
}
\]

%%%%%%%%%%%%%%%%%%%%%%%%%%%%%%%%%%%%%%%%%%%%%%%%%%%%%%%%%%%%
\subsection{Worked Example: Logistic Model}
\label{subsec:worked-logistic-model}
%%%%%%%%%%%%%%%%%%%%%%%%%%%%%%%%%%%%%%%%%%%%%%%%%%%%%%%%%%%%

\begin{workedexamplebox}{Carrying Capacity}

Suppose

\[
P'
=
0.4P
\left(
1-\frac{P}{1000}
\right),
\qquad
P(0)=100.
\]

Here,

\[
K=1000,
\qquad
r=0.4.
\]

The constant is

\[
A
=
\frac{
1000-100
}{
100
}
=
9.
\]

Therefore,

\begin{answerbox}
\[
\boxed{
P(t)
=
\frac{
1000
}{
1+9e^{-0.4t}
}.
}
\]
\end{answerbox}

\end{workedexamplebox}

%%%%%%%%%%%%%%%%%%%%%%%%%%%%%%%%%%%%%%%%%%%%%%%%%%%%%%%%%%%%
\subsection{Newton's Law of Cooling}
\label{subsec:newtons-law-cooling}
%%%%%%%%%%%%%%%%%%%%%%%%%%%%%%%%%%%%%%%%%%%%%%%%%%%%%%%%%%%%

Newton's law of cooling assumes the rate of temperature change is
proportional to the difference between an object's temperature and the
ambient temperature.

Let \(T_a\) be constant ambient temperature.

Then

\[
\boxed{
\frac{dT}{dt}
=
-k
(
T-T_a
),
\qquad
k>0.
}
\]

The solution is

\[
\boxed{
T(t)
=
T_a
+
(
T_0-T_a
)
e^{-kt}.
}
\]

%%%%%%%%%%%%%%%%%%%%%%%%%%%%%%%%%%%%%%%%%%%%%%%%%%%%%%%%%%%%
\subsection{Mixing Problems}
\label{subsec:mixing-problems}
%%%%%%%%%%%%%%%%%%%%%%%%%%%%%%%%%%%%%%%%%%%%%%%%%%%%%%%%%%%%

Let

\[
A(t)
\]

be the amount of a substance in a well-mixed tank.

The general balance law is

\[
\boxed{
\frac{dA}{dt}
=
\text{rate in}
-
\text{rate out}.
}
\]

If concentration is uniform,

\[
\boxed{
\text{concentration}
=
\frac{
A(t)
}{
V(t)
}.
}
\]

This determines the outgoing substance rate.

%%%%%%%%%%%%%%%%%%%%%%%%%%%%%%%%%%%%%%%%%%%%%%%%%%%%%%%%%%%%
\subsection{Second-Order Linear Equations}
\label{subsec:second-order-linear-odes}
%%%%%%%%%%%%%%%%%%%%%%%%%%%%%%%%%%%%%%%%%%%%%%%%%%%%%%%%%%%%

A second-order linear equation has form

\[
\boxed{
a(x)y''
+
b(x)y'
+
c(x)y
=
g(x).
}
\]

When \(a,b,c\) are constants, characteristic-equation methods become
available.

%%%%%%%%%%%%%%%%%%%%%%%%%%%%%%%%%%%%%%%%%%%%%%%%%%%%%%%%%%%%
\subsection{Superposition Principle}
\label{subsec:ode-superposition}
%%%%%%%%%%%%%%%%%%%%%%%%%%%%%%%%%%%%%%%%%%%%%%%%%%%%%%%%%%%%

\begin{theorem}[Superposition Principle]
If \(y_1\) and \(y_2\) solve the homogeneous linear equation

\[
L[y]=0,
\]

then

\[
\boxed{
C_1y_1+C_2y_2
}
\]

also solves the equation for arbitrary constants \(C_1,C_2\).
\end{theorem}

This follows from linearity of the differential operator \(L\).

%%%%%%%%%%%%%%%%%%%%%%%%%%%%%%%%%%%%%%%%%%%%%%%%%%%%%%%%%%%%
\subsection{Linear Independence of Solutions}
\label{subsec:linear-independence-ode-solutions}
%%%%%%%%%%%%%%%%%%%%%%%%%%%%%%%%%%%%%%%%%%%%%%%%%%%%%%%%%%%%

Two functions \(y_1,y_2\) are linearly independent on an interval if

\[
\boxed{
C_1y_1+C_2y_2=0
}
\]

throughout the interval implies

\[
\boxed{
C_1=C_2=0.
}
\]

Two linearly independent solutions form a fundamental solution set for a
second-order homogeneous linear equation.

%%%%%%%%%%%%%%%%%%%%%%%%%%%%%%%%%%%%%%%%%%%%%%%%%%%%%%%%%%%%
\subsection{Wronskian}
\label{subsec:wronskian}
%%%%%%%%%%%%%%%%%%%%%%%%%%%%%%%%%%%%%%%%%%%%%%%%%%%%%%%%%%%%

For functions \(y_1,y_2\), the Wronskian is

\[
\boxed{
W(y_1,y_2)
=
\begin{vmatrix}
y_1 & y_2\\
y_1' & y_2'
\end{vmatrix}
=
y_1y_2'
-
y_2y_1'.
}
\]

If

\[
W(x_0)\neq0
\]

at some point under the usual linear-ODE hypotheses, the solutions are
linearly independent.

%%%%%%%%%%%%%%%%%%%%%%%%%%%%%%%%%%%%%%%%%%%%%%%%%%%%%%%%%%%%
\subsection{Abel's Identity Preview}
\label{subsec:abel-identity-preview}
%%%%%%%%%%%%%%%%%%%%%%%%%%%%%%%%%%%%%%%%%%%%%%%%%%%%%%%%%%%%

For

\[
y''
+
P(x)y'
+
Q(x)y
=
0,
\]

the Wronskian satisfies

\[
\boxed{
W(x)
=
W(x_0)
\exp
\left[
-
\int_{x_0}^{x}
P(s)\,ds
\right].
}
\]

Thus if the Wronskian is nonzero at one point, it remains nonzero on the
interval.

%%%%%%%%%%%%%%%%%%%%%%%%%%%%%%%%%%%%%%%%%%%%%%%%%%%%%%%%%%%%
\subsection{Constant-Coefficient Homogeneous Equations}
\label{subsec:constant-coefficient-homogeneous}
%%%%%%%%%%%%%%%%%%%%%%%%%%%%%%%%%%%%%%%%%%%%%%%%%%%%%%%%%%%%

Consider

\[
\boxed{
ay''
+
by'
+
cy
=
0,
}
\]

where \(a\neq0\).

Try

\[
\boxed{
y=e^{rx}.
}
\]

Then

\[
y'=re^{rx},
\qquad
y''=r^2e^{rx}.
\]

Substitution gives

\[
\boxed{
ar^2+br+c=0.
}
\]

This is the characteristic equation.

%%%%%%%%%%%%%%%%%%%%%%%%%%%%%%%%%%%%%%%%%%%%%%%%%%%%%%%%%%%%
\subsection{Distinct Real Characteristic Roots}
\label{subsec:distinct-real-roots}
%%%%%%%%%%%%%%%%%%%%%%%%%%%%%%%%%%%%%%%%%%%%%%%%%%%%%%%%%%%%

If the characteristic equation has distinct real roots

\[
r_1\neq r_2,
\]

then

\[
\boxed{
y
=
C_1e^{r_1x}
+
C_2e^{r_2x}.
}
\]

%%%%%%%%%%%%%%%%%%%%%%%%%%%%%%%%%%%%%%%%%%%%%%%%%%%%%%%%%%%%
\subsection{Worked Example: Distinct Roots}
\label{subsec:worked-distinct-roots}
%%%%%%%%%%%%%%%%%%%%%%%%%%%%%%%%%%%%%%%%%%%%%%%%%%%%%%%%%%%%

\begin{workedexamplebox}{Second-Order Homogeneous Equation}

Solve

\[
y''-5y'+6y=0.
\]

The characteristic equation is

\[
r^2-5r+6=0.
\]

Factor:

\[
(r-2)(r-3)=0.
\]

Thus

\[
r_1=2,
\qquad
r_2=3.
\]

Therefore,

\begin{answerbox}
\[
\boxed{
y
=
C_1e^{2x}
+
C_2e^{3x}.
}
\]
\end{answerbox}

\end{workedexamplebox}

%%%%%%%%%%%%%%%%%%%%%%%%%%%%%%%%%%%%%%%%%%%%%%%%%%%%%%%%%%%%
\subsection{Repeated Real Root}
\label{subsec:repeated-characteristic-root}
%%%%%%%%%%%%%%%%%%%%%%%%%%%%%%%%%%%%%%%%%%%%%%%%%%%%%%%%%%%%

If the characteristic equation has repeated root

\[
r,
\]

then one solution is

\[
e^{rx}.
\]

A second independent solution is

\[
xe^{rx}.
\]

Thus,

\[
\boxed{
y
=
(
C_1+C_2x
)
e^{rx}.
}
\]

%%%%%%%%%%%%%%%%%%%%%%%%%%%%%%%%%%%%%%%%%%%%%%%%%%%%%%%%%%%%
\subsection{Worked Example: Repeated Root}
\label{subsec:worked-repeated-root}
%%%%%%%%%%%%%%%%%%%%%%%%%%%%%%%%%%%%%%%%%%%%%%%%%%%%%%%%%%%%

\begin{workedexamplebox}{Repeated Characteristic Root}

Solve

\[
y''-4y'+4y=0.
\]

The characteristic equation is

\[
r^2-4r+4=0,
\]

so

\[
(r-2)^2=0.
\]

Therefore,

\begin{answerbox}
\[
\boxed{
y
=
(
C_1+C_2x
)
e^{2x}.
}
\]
\end{answerbox}

\end{workedexamplebox}

%%%%%%%%%%%%%%%%%%%%%%%%%%%%%%%%%%%%%%%%%%%%%%%%%%%%%%%%%%%%
\subsection{Complex Characteristic Roots}
\label{subsec:complex-characteristic-roots}
%%%%%%%%%%%%%%%%%%%%%%%%%%%%%%%%%%%%%%%%%%%%%%%%%%%%%%%%%%%%

Suppose the roots are

\[
\boxed{
r
=
\alpha
\pm
i\beta.
}
\]

Then a real fundamental solution pair is

\[
e^{\alpha x}\cos(\beta x),
\]

and

\[
e^{\alpha x}\sin(\beta x).
\]

Therefore,

\[
\boxed{
y
=
e^{\alpha x}
\left[
C_1\cos(\beta x)
+
C_2\sin(\beta x)
\right].
}
\]

%%%%%%%%%%%%%%%%%%%%%%%%%%%%%%%%%%%%%%%%%%%%%%%%%%%%%%%%%%%%
\subsection{Worked Example: Complex Roots}
\label{subsec:worked-complex-roots}
%%%%%%%%%%%%%%%%%%%%%%%%%%%%%%%%%%%%%%%%%%%%%%%%%%%%%%%%%%%%

\begin{workedexamplebox}{Oscillatory Solution}

Solve

\[
y''+4y=0.
\]

The characteristic equation is

\[
r^2+4=0.
\]

Thus,

\[
r=\pm2i.
\]

Therefore,

\begin{answerbox}
\[
\boxed{
y
=
C_1\cos(2x)
+
C_2\sin(2x).
}
\]
\end{answerbox}

\end{workedexamplebox}

%%%%%%%%%%%%%%%%%%%%%%%%%%%%%%%%%%%%%%%%%%%%%%%%%%%%%%%%%%%%
\subsection{Higher-Order Constant-Coefficient Equations}
\label{subsec:higher-order-constant-coefficient}
%%%%%%%%%%%%%%%%%%%%%%%%%%%%%%%%%%%%%%%%%%%%%%%%%%%%%%%%%%%%

For

\[
\boxed{
a_ny^{(n)}
+
\cdots
+
a_1y'
+
a_0y
=
0,
}
\]

the characteristic polynomial is

\[
\boxed{
a_nr^n
+
\cdots
+
a_1r
+
a_0
=
0.
}
\]

Each root contributes solutions according to its multiplicity.

If \(r\) has multiplicity \(m\), then

\[
\boxed{
e^{rx},
xe^{rx},
\ldots,
x^{m-1}e^{rx}
}
\]

appear.

%%%%%%%%%%%%%%%%%%%%%%%%%%%%%%%%%%%%%%%%%%%%%%%%%%%%%%%%%%%%
\subsection{Nonhomogeneous Linear Equations}
\label{subsec:linear-nonhomogeneous-solution}
%%%%%%%%%%%%%%%%%%%%%%%%%%%%%%%%%%%%%%%%%%%%%%%%%%%%%%%%%%%%

For

\[
\boxed{
L[y]
=
g(x),
}
\]

the general solution is

\[
\boxed{
y
=
y_h+y_p,
}
\]

where

\[
y_h
\]

solves

\[
L[y_h]=0,
\]

and \(y_p\) is one particular solution of

\[
L[y_p]=g(x).
\]

%%%%%%%%%%%%%%%%%%%%%%%%%%%%%%%%%%%%%%%%%%%%%%%%%%%%%%%%%%%%
\subsection{Why \(y_h+y_p\) Works}
\label{subsec:why-homogeneous-plus-particular}
%%%%%%%%%%%%%%%%%%%%%%%%%%%%%%%%%%%%%%%%%%%%%%%%%%%%%%%%%%%%

By linearity,

\[
L[y_h+y_p]
=
L[y_h]
+
L[y_p].
\]

Since

\[
L[y_h]=0
\]

and

\[
L[y_p]=g,
\]

we obtain

\[
\boxed{
L[y_h+y_p]
=
g.
}
\]

%%%%%%%%%%%%%%%%%%%%%%%%%%%%%%%%%%%%%%%%%%%%%%%%%%%%%%%%%%%%
\subsection{Undetermined Coefficients}
\label{subsec:undetermined-coefficients}
%%%%%%%%%%%%%%%%%%%%%%%%%%%%%%%%%%%%%%%%%%%%%%%%%%%%%%%%%%%%

The method of undetermined coefficients applies to constant-coefficient
linear equations when the forcing term belongs to certain families,
including:

\[
\boxed{
\text{polynomials},
}
\]

\[
\boxed{
e^{ax},
}
\]

\[
\boxed{
\sin bx,
}
\]

\[
\boxed{
\cos bx,
}
\]

and finite products or sums of these.

%%%%%%%%%%%%%%%%%%%%%%%%%%%%%%%%%%%%%%%%%%%%%%%%%%%%%%%%%%%%
\subsection{Trial Forms}
\label{subsec:undetermined-coefficient-trials}
%%%%%%%%%%%%%%%%%%%%%%%%%%%%%%%%%%%%%%%%%%%%%%%%%%%%%%%%%%%%

For forcing

\[
g(x)=Ae^{ax},
\]

try

\[
\boxed{
y_p=Ce^{ax}.
}
\]

For

\[
g(x)
=
P_n(x),
\]

try a general polynomial of the same degree.

For

\[
g(x)
=
A\cos bx+B\sin bx,
\]

try

\[
\boxed{
y_p
=
C\cos bx
+
D\sin bx.
}
\]

%%%%%%%%%%%%%%%%%%%%%%%%%%%%%%%%%%%%%%%%%%%%%%%%%%%%%%%%%%%%
\subsection{Resonance in Undetermined Coefficients}
\label{subsec:undetermined-coefficients-resonance}
%%%%%%%%%%%%%%%%%%%%%%%%%%%%%%%%%%%%%%%%%%%%%%%%%%%%%%%%%%%%

If the proposed trial already appears in the homogeneous solution, it
must be multiplied by a sufficient power of \(x\).

For example, if

\[
e^{ax}
\]

is already a homogeneous solution of multiplicity \(m\), use

\[
\boxed{
x^mCe^{ax}.
}
\]

%%%%%%%%%%%%%%%%%%%%%%%%%%%%%%%%%%%%%%%%%%%%%%%%%%%%%%%%%%%%
\subsection{Worked Example: Undetermined Coefficients}
\label{subsec:worked-undetermined-coefficients}
%%%%%%%%%%%%%%%%%%%%%%%%%%%%%%%%%%%%%%%%%%%%%%%%%%%%%%%%%%%%

\begin{workedexamplebox}{Constant Forcing}

Solve

\[
y''-y=2.
\]

The homogeneous equation is

\[
y''-y=0.
\]

Its characteristic equation is

\[
r^2-1=0,
\]

giving

\[
y_h
=
C_1e^x
+
C_2e^{-x}.
\]

For constant forcing \(2\), try

\[
y_p=A.
\]

Then

\[
y_p''=0.
\]

Substitute:

\[
0-A=2.
\]

Thus

\[
A=-2.
\]

Therefore,

\begin{answerbox}
\[
\boxed{
y
=
C_1e^x
+
C_2e^{-x}
-
2.
}
\]
\end{answerbox}

\end{workedexamplebox}

%%%%%%%%%%%%%%%%%%%%%%%%%%%%%%%%%%%%%%%%%%%%%%%%%%%%%%%%%%%%
\subsection{Variation of Parameters}
\label{subsec:variation-of-parameters}
%%%%%%%%%%%%%%%%%%%%%%%%%%%%%%%%%%%%%%%%%%%%%%%%%%%%%%%%%%%%

Variation of parameters finds a particular solution of

\[
\boxed{
y''
+
P(x)y'
+
Q(x)y
=
g(x)
}
\]

using a fundamental homogeneous pair

\[
y_1,
\qquad
y_2.
\]

Assume

\[
\boxed{
y_p
=
u_1(x)y_1(x)
+
u_2(x)y_2(x).
}
\]

%%%%%%%%%%%%%%%%%%%%%%%%%%%%%%%%%%%%%%%%%%%%%%%%%%%%%%%%%%%%
\subsection{Variation-of-Parameters Formulas}
\label{subsec:variation-parameters-formulas}
%%%%%%%%%%%%%%%%%%%%%%%%%%%%%%%%%%%%%%%%%%%%%%%%%%%%%%%%%%%%

With Wronskian

\[
W
=
y_1y_2'
-
y_2y_1',
\]

one obtains

\[
\boxed{
u_1'
=
-
\frac{
y_2g
}{
W
},
}
\]

and

\[
\boxed{
u_2'
=
\frac{
y_1g
}{
W
}.
}
\]

Therefore,

\[
\boxed{
u_1
=
-
\int
\frac{
y_2g
}{
W
}
\,dx,
}
\]

\[
\boxed{
u_2
=
\int
\frac{
y_1g
}{
W
}
\,dx.
}
\]

%%%%%%%%%%%%%%%%%%%%%%%%%%%%%%%%%%%%%%%%%%%%%%%%%%%%%%%%%%%%
\subsection{Advantage of Variation of Parameters}
\label{subsec:variation-parameters-advantage}
%%%%%%%%%%%%%%%%%%%%%%%%%%%%%%%%%%%%%%%%%%%%%%%%%%%%%%%%%%%%

Unlike undetermined coefficients, variation of parameters does not
require the forcing function to have a special polynomial-exponential-
trigonometric form.

Its main requirements are knowledge of a fundamental homogeneous
solution set and integrability of the resulting expressions.

%%%%%%%%%%%%%%%%%%%%%%%%%%%%%%%%%%%%%%%%%%%%%%%%%%%%%%%%%%%%
\subsection{Reduction of Order}
\label{subsec:reduction-of-order}
%%%%%%%%%%%%%%%%%%%%%%%%%%%%%%%%%%%%%%%%%%%%%%%%%%%%%%%%%%%%

Suppose one solution

\[
y_1(x)
\]

of a second-order homogeneous linear equation is known.

A second solution may be sought in the form

\[
\boxed{
y_2(x)
=
v(x)y_1(x).
}
\]

Substituting reduces the order of the unknown equation for \(v\).

%%%%%%%%%%%%%%%%%%%%%%%%%%%%%%%%%%%%%%%%%%%%%%%%%%%%%%%%%%%%
\subsection{Reduction-of-Order Formula}
\label{subsec:reduction-order-formula}
%%%%%%%%%%%%%%%%%%%%%%%%%%%%%%%%%%%%%%%%%%%%%%%%%%%%%%%%%%%%

For

\[
y''
+
P(x)y'
+
Q(x)y
=
0,
\]

if \(y_1\neq0\) is known, a second solution can be written

\[
\boxed{
y_2
=
y_1
\int
\frac{
e^{-\int P(x)\,dx}
}{
y_1^2
}
\,dx.
}
\]

This is useful when only one solution is easily recognized.

%%%%%%%%%%%%%%%%%%%%%%%%%%%%%%%%%%%%%%%%%%%%%%%%%%%%%%%%%%%%
\subsection{Mechanical Vibrations}
\label{subsec:mechanical-vibrations}
%%%%%%%%%%%%%%%%%%%%%%%%%%%%%%%%%%%%%%%%%%%%%%%%%%%%%%%%%%%%

A spring--mass--damper system is modeled by

\[
\boxed{
m x''
+
c x'
+
k x
=
F(t),
}
\]

where:

\[
m
=
\text{mass},
\]

\[
c
=
\text{damping coefficient},
\]

\[
k
=
\text{spring constant},
\]

and

\[
F(t)
=
\text{external forcing}.
\]

%%%%%%%%%%%%%%%%%%%%%%%%%%%%%%%%%%%%%%%%%%%%%%%%%%%%%%%%%%%%
\subsection{Free Undamped Motion}
\label{subsec:free-undamped-motion}
%%%%%%%%%%%%%%%%%%%%%%%%%%%%%%%%%%%%%%%%%%%%%%%%%%%%%%%%%%%%

If

\[
c=0
\]

and

\[
F(t)=0,
\]

then

\[
\boxed{
mx''+kx=0.
}
\]

Divide by \(m\):

\[
x''
+
\omega_0^2x
=
0,
\]

where

\[
\boxed{
\omega_0
=
\sqrt{
\frac{k}{m}
}
}
\]

is the natural angular frequency.

%%%%%%%%%%%%%%%%%%%%%%%%%%%%%%%%%%%%%%%%%%%%%%%%%%%%%%%%%%%%
\subsection{Simple Harmonic Motion}
\label{subsec:simple-harmonic-motion}
%%%%%%%%%%%%%%%%%%%%%%%%%%%%%%%%%%%%%%%%%%%%%%%%%%%%%%%%%%%%

The solution is

\[
\boxed{
x(t)
=
C_1\cos(\omega_0t)
+
C_2\sin(\omega_0t).
}
\]

Equivalently,

\[
\boxed{
x(t)
=
A
\cos
(
\omega_0t-\phi
).
}
\]

The period is

\[
\boxed{
T
=
\frac{
2\pi
}{
\omega_0
}.
}
\]

%%%%%%%%%%%%%%%%%%%%%%%%%%%%%%%%%%%%%%%%%%%%%%%%%%%%%%%%%%%%
\subsection{Damped Motion}
\label{subsec:damped-motion}
%%%%%%%%%%%%%%%%%%%%%%%%%%%%%%%%%%%%%%%%%%%%%%%%%%%%%%%%%%%%

For

\[
mx''
+
cx'
+
kx
=
0,
\]

the characteristic equation is

\[
\boxed{
mr^2+cr+k=0.
}
\]

The discriminant

\[
\boxed{
c^2-4mk
}
\]

determines the type of damping.

%%%%%%%%%%%%%%%%%%%%%%%%%%%%%%%%%%%%%%%%%%%%%%%%%%%%%%%%%%%%
\subsection{Overdamping}
\label{subsec:overdamping}
%%%%%%%%%%%%%%%%%%%%%%%%%%%%%%%%%%%%%%%%%%%%%%%%%%%%%%%%%%%%

If

\[
\boxed{
c^2>4mk,
}
\]

there are two distinct negative real roots under ordinary physical
parameters.

The system returns to equilibrium without oscillating.

This is overdamped motion.

%%%%%%%%%%%%%%%%%%%%%%%%%%%%%%%%%%%%%%%%%%%%%%%%%%%%%%%%%%%%
\subsection{Critical Damping}
\label{subsec:critical-damping}
%%%%%%%%%%%%%%%%%%%%%%%%%%%%%%%%%%%%%%%%%%%%%%%%%%%%%%%%%%%%

If

\[
\boxed{
c^2=4mk,
}
\]

the characteristic equation has a repeated root.

The system returns to equilibrium as rapidly as possible without
oscillation in the standard linear model.

This is critical damping.

%%%%%%%%%%%%%%%%%%%%%%%%%%%%%%%%%%%%%%%%%%%%%%%%%%%%%%%%%%%%
\subsection{Underdamping}
\label{subsec:underdamping}
%%%%%%%%%%%%%%%%%%%%%%%%%%%%%%%%%%%%%%%%%%%%%%%%%%%%%%%%%%%%

If

\[
\boxed{
c^2<4mk,
}
\]

the roots are complex:

\[
\boxed{
r
=
-\frac{c}{2m}
\pm
i\omega_d.
}
\]

The displacement oscillates with exponentially decreasing amplitude.

This is underdamped motion.

%%%%%%%%%%%%%%%%%%%%%%%%%%%%%%%%%%%%%%%%%%%%%%%%%%%%%%%%%%%%
\subsection{Forced Oscillations}
\label{subsec:forced-oscillations}
%%%%%%%%%%%%%%%%%%%%%%%%%%%%%%%%%%%%%%%%%%%%%%%%%%%%%%%%%%%%

With periodic forcing,

\[
\boxed{
mx''
+
cx'
+
kx
=
F_0\cos(\omega t).
}
\]

The total motion contains:

\[
\boxed{
\text{transient response}
}
\]

and

\[
\boxed{
\text{steady-state response}.
}
\]

The transient usually decays when damping is present.

%%%%%%%%%%%%%%%%%%%%%%%%%%%%%%%%%%%%%%%%%%%%%%%%%%%%%%%%%%%%
\subsection{Resonance}
\label{subsec:ode-resonance}
%%%%%%%%%%%%%%%%%%%%%%%%%%%%%%%%%%%%%%%%%%%%%%%%%%%%%%%%%%%%

In an undamped system, forcing near the natural frequency can create very
large oscillations.

For the idealized equation

\[
\boxed{
x''
+
\omega_0^2x
=
F_0\cos(\omega_0t),
}
\]

the forcing resonates with the homogeneous solution.

A particular solution contains a factor of \(t\), producing growing
amplitude.

%%%%%%%%%%%%%%%%%%%%%%%%%%%%%%%%%%%%%%%%%%%%%%%%%%%%%%%%%%%%
\subsection{Systems of First-Order ODEs}
\label{subsec:first-order-ode-systems}
%%%%%%%%%%%%%%%%%%%%%%%%%%%%%%%%%%%%%%%%%%%%%%%%%%%%%%%%%%%%

A system may be written

\[
\boxed{
\mathbf{x}'
=
\mathbf{f}
(
t,\mathbf{x}
).
}
\]

For two variables,

\[
\boxed{
\begin{cases}
x'=f(t,x,y),\\
y'=g(t,x,y).
\end{cases}
}
\]

Systems allow interacting quantities to evolve simultaneously.

%%%%%%%%%%%%%%%%%%%%%%%%%%%%%%%%%%%%%%%%%%%%%%%%%%%%%%%%%%%%
\subsection{Higher-Order ODE as a First-Order System}
\label{subsec:higher-order-to-system}
%%%%%%%%%%%%%%%%%%%%%%%%%%%%%%%%%%%%%%%%%%%%%%%%%%%%%%%%%%%%

Consider

\[
y''=F(t,y,y').
\]

Define

\[
\boxed{
x_1=y,
\qquad
x_2=y'.
}
\]

Then

\[
\boxed{
x_1'=x_2,
}
\]

and

\[
\boxed{
x_2'
=
F(t,x_1,x_2).
}
\]

Thus every higher-order ODE can be rewritten as a first-order system by
introducing derivative variables.

%%%%%%%%%%%%%%%%%%%%%%%%%%%%%%%%%%%%%%%%%%%%%%%%%%%%%%%%%%%%
\subsection{Linear Systems}
\label{subsec:linear-ode-systems}
%%%%%%%%%%%%%%%%%%%%%%%%%%%%%%%%%%%%%%%%%%%%%%%%%%%%%%%%%%%%

A linear autonomous system has form

\[
\boxed{
\mathbf{x}'
=
A\mathbf{x}.
}
\]

A nonhomogeneous linear system is

\[
\boxed{
\mathbf{x}'
=
A(t)\mathbf{x}
+
\mathbf{g}(t).
}
\]

Matrix methods provide a natural solution framework.

%%%%%%%%%%%%%%%%%%%%%%%%%%%%%%%%%%%%%%%%%%%%%%%%%%%%%%%%%%%%
\subsection{Eigenvalue Method}
\label{subsec:eigenvalue-method-ode}
%%%%%%%%%%%%%%%%%%%%%%%%%%%%%%%%%%%%%%%%%%%%%%%%%%%%%%%%%%%%

For

\[
\mathbf{x}'
=
A\mathbf{x},
\]

seek a solution

\[
\boxed{
\mathbf{x}(t)
=
\mathbf{v}e^{\lambda t}.
}
\]

Then

\[
\mathbf{x}'
=
\lambda
\mathbf{v}e^{\lambda t}.
\]

Substituting gives

\[
\lambda
\mathbf{v}
=
A\mathbf{v}.
\]

Thus \(\lambda\) and \(\mathbf{v}\) must be an eigenvalue-eigenvector pair
of \(A\).

%%%%%%%%%%%%%%%%%%%%%%%%%%%%%%%%%%%%%%%%%%%%%%%%%%%%%%%%%%%%
\subsection{Distinct Eigenvalues}
\label{subsec:distinct-eigenvalue-system}
%%%%%%%%%%%%%%%%%%%%%%%%%%%%%%%%%%%%%%%%%%%%%%%%%%%%%%%%%%%%

If \(A\) has \(n\) linearly independent eigenvectors

\[
\mathbf{v}_1,
\ldots,
\mathbf{v}_n
\]

with eigenvalues

\[
\lambda_1,
\ldots,
\lambda_n,
\]

then

\[
\boxed{
\mathbf{x}(t)
=
\sum_{j=1}^{n}
C_j
\mathbf{v}_j
e^{\lambda_jt}.
}
\]

%%%%%%%%%%%%%%%%%%%%%%%%%%%%%%%%%%%%%%%%%%%%%%%%%%%%%%%%%%%%
\subsection{Worked Example: Linear System}
\label{subsec:worked-linear-system}
%%%%%%%%%%%%%%%%%%%%%%%%%%%%%%%%%%%%%%%%%%%%%%%%%%%%%%%%%%%%

\begin{workedexamplebox}{Diagonal System}

Solve

\[
\mathbf{x}'
=
\begin{pmatrix}
2&0\\
0&-1
\end{pmatrix}
\mathbf{x}.
\]

The eigenvalues are

\[
\lambda_1=2,
\qquad
\lambda_2=-1,
\]

with eigenvectors

\[
\mathbf{v}_1
=
\begin{pmatrix}
1\\
0
\end{pmatrix},
\qquad
\mathbf{v}_2
=
\begin{pmatrix}
0\\
1
\end{pmatrix}.
\]

Therefore,

\begin{answerbox}
\[
\boxed{
\mathbf{x}(t)
=
C_1
\begin{pmatrix}
1\\
0
\end{pmatrix}
e^{2t}
+
C_2
\begin{pmatrix}
0\\
1
\end{pmatrix}
e^{-t}.
}
\]
\end{answerbox}

\end{workedexamplebox}

%%%%%%%%%%%%%%%%%%%%%%%%%%%%%%%%%%%%%%%%%%%%%%%%%%%%%%%%%%%%
\subsection{Matrix Exponential}
\label{subsec:matrix-exponential-ode}
%%%%%%%%%%%%%%%%%%%%%%%%%%%%%%%%%%%%%%%%%%%%%%%%%%%%%%%%%%%%

The solution of

\[
\mathbf{x}'
=
A\mathbf{x},
\qquad
\mathbf{x}(0)=\mathbf{x}_0,
\]

can be written

\[
\boxed{
\mathbf{x}(t)
=
e^{At}
\mathbf{x}_0.
}
\]

The matrix exponential is

\[
\boxed{
e^{At}
=
I
+
At
+
\frac{
A^2t^2
}{
2!
}
+
\frac{
A^3t^3
}{
3!
}
+
\cdots.
}
\]

%%%%%%%%%%%%%%%%%%%%%%%%%%%%%%%%%%%%%%%%%%%%%%%%%%%%%%%%%%%%
\subsection{Fundamental Matrix}
\label{subsec:fundamental-matrix-ode}
%%%%%%%%%%%%%%%%%%%%%%%%%%%%%%%%%%%%%%%%%%%%%%%%%%%%%%%%%%%%

A matrix \(\Phi(t)\) whose columns form a fundamental solution set of

\[
\mathbf{x}'
=
A(t)\mathbf{x}
\]

is called a fundamental matrix.

It satisfies

\[
\boxed{
\Phi'(t)
=
A(t)\Phi(t).
}
\]

For a constant matrix \(A\),

\[
\boxed{
\Phi(t)=e^{At}
}
\]

is a fundamental matrix.

%%%%%%%%%%%%%%%%%%%%%%%%%%%%%%%%%%%%%%%%%%%%%%%%%%%%%%%%%%%%
\subsection{Planar Linear Systems}
\label{subsec:planar-linear-systems}
%%%%%%%%%%%%%%%%%%%%%%%%%%%%%%%%%%%%%%%%%%%%%%%%%%%%%%%%%%%%

For

\[
\boxed{
\mathbf{x}'
=
A\mathbf{x},
\qquad
A\in\mathbb{R}^{2\times2},
}
\]

the equilibrium

\[
\mathbf{x}=\mathbf{0}
\]

can often be classified using eigenvalues.

The phase portrait shows trajectories in the \(xy\)-plane.

%%%%%%%%%%%%%%%%%%%%%%%%%%%%%%%%%%%%%%%%%%%%%%%%%%%%%%%%%%%%
\subsection{Stable Node}
\label{subsec:stable-node}
%%%%%%%%%%%%%%%%%%%%%%%%%%%%%%%%%%%%%%%%%%%%%%%%%%%%%%%%%%%%

If both eigenvalues are real, distinct, and negative,

\[
\boxed{
\lambda_1<0,
\qquad
\lambda_2<0,
}
\]

then the origin is a stable node.

Nearby trajectories approach the origin.

%%%%%%%%%%%%%%%%%%%%%%%%%%%%%%%%%%%%%%%%%%%%%%%%%%%%%%%%%%%%
\subsection{Unstable Node}
\label{subsec:unstable-node}
%%%%%%%%%%%%%%%%%%%%%%%%%%%%%%%%%%%%%%%%%%%%%%%%%%%%%%%%%%%%

If both eigenvalues are real and positive,

\[
\boxed{
\lambda_1>0,
\qquad
\lambda_2>0,
}
\]

the origin is an unstable node.

Nearby trajectories move away.

%%%%%%%%%%%%%%%%%%%%%%%%%%%%%%%%%%%%%%%%%%%%%%%%%%%%%%%%%%%%
\subsection{Saddle Point}
\label{subsec:saddle-point-linear-system}
%%%%%%%%%%%%%%%%%%%%%%%%%%%%%%%%%%%%%%%%%%%%%%%%%%%%%%%%%%%%

If the eigenvalues have opposite signs,

\[
\boxed{
\lambda_1\lambda_2<0,
}
\]

the origin is a saddle point.

One eigendirection is attracting and the other is repelling.

A saddle point is unstable.

%%%%%%%%%%%%%%%%%%%%%%%%%%%%%%%%%%%%%%%%%%%%%%%%%%%%%%%%%%%%
\subsection{Spiral Point}
\label{subsec:spiral-point-linear-system}
%%%%%%%%%%%%%%%%%%%%%%%%%%%%%%%%%%%%%%%%%%%%%%%%%%%%%%%%%%%%

If

\[
\boxed{
\lambda
=
\alpha
\pm
i\beta,
\qquad
\beta\neq0,
}
\]

solutions rotate while their amplitude behaves approximately like

\[
e^{\alpha t}.
\]

If

\[
\boxed{
\alpha<0,
}
\]

the spiral is stable.

If

\[
\boxed{
\alpha>0,
}
\]

it is unstable.

%%%%%%%%%%%%%%%%%%%%%%%%%%%%%%%%%%%%%%%%%%%%%%%%%%%%%%%%%%%%
\subsection{Center}
\label{subsec:center-linear-system}
%%%%%%%%%%%%%%%%%%%%%%%%%%%%%%%%%%%%%%%%%%%%%%%%%%%%%%%%%%%%

If the eigenvalues are purely imaginary,

\[
\boxed{
\lambda
=
\pm i\beta,
}
\]

a linear planar system may have closed trajectories surrounding the
origin.

The equilibrium is a center.

It is stable in the Lyapunov sense for the ideal linear system but not
asymptotically stable.

%%%%%%%%%%%%%%%%%%%%%%%%%%%%%%%%%%%%%%%%%%%%%%%%%%%%%%%%%%%%
\subsection{Trace--Determinant Classification}
\label{subsec:trace-determinant-classification}
%%%%%%%%%%%%%%%%%%%%%%%%%%%%%%%%%%%%%%%%%%%%%%%%%%%%%%%%%%%%

For

\[
A
=
\begin{pmatrix}
a&b\\
c&d
\end{pmatrix},
\]

define

\[
\boxed{
\tau
=
\operatorname{tr}(A)
=
a+d,
}
\]

and

\[
\boxed{
\Delta
=
\det(A)
=
ad-bc.
}
\]

The characteristic polynomial is

\[
\boxed{
\lambda^2
-
\tau\lambda
+
\Delta
=
0.
}
\]

%%%%%%%%%%%%%%%%%%%%%%%%%%%%%%%%%%%%%%%%%%%%%%%%%%%%%%%%%%%%
\subsection{Trace--Determinant Discriminant}
\label{subsec:trace-determinant-discriminant}
%%%%%%%%%%%%%%%%%%%%%%%%%%%%%%%%%%%%%%%%%%%%%%%%%%%%%%%%%%%%

The eigenvalue discriminant is

\[
\boxed{
D
=
\tau^2
-
4\Delta.
}
\]

Then:

\[
\boxed{
D>0
\rightarrow
\text{real eigenvalues},
}
\]

\[
\boxed{
D<0
\rightarrow
\text{complex-conjugate eigenvalues}.
}
\]

If

\[
\Delta<0,
\]

the equilibrium is a saddle.

%%%%%%%%%%%%%%%%%%%%%%%%%%%%%%%%%%%%%%%%%%%%%%%%%%%%%%%%%%%%
\subsection{Nonlinear Systems}
\label{subsec:nonlinear-ode-systems}
%%%%%%%%%%%%%%%%%%%%%%%%%%%%%%%%%%%%%%%%%%%%%%%%%%%%%%%%%%%%

A nonlinear autonomous system has form

\[
\boxed{
\mathbf{x}'
=
\mathbf{f}(\mathbf{x}).
}
\]

Equilibria satisfy

\[
\boxed{
\mathbf{f}(\mathbf{x}^*)
=
\mathbf{0}.
}
\]

Exact solutions are often unavailable.

Qualitative analysis becomes central.

%%%%%%%%%%%%%%%%%%%%%%%%%%%%%%%%%%%%%%%%%%%%%%%%%%%%%%%%%%%%
\subsection{Linearization}
\label{subsec:ode-linearization}
%%%%%%%%%%%%%%%%%%%%%%%%%%%%%%%%%%%%%%%%%%%%%%%%%%%%%%%%%%%%

Near an equilibrium \(\mathbf{x}^*\), approximate

\[
\mathbf{f}(\mathbf{x})
\]

using the Jacobian matrix

\[
\boxed{
J
=
D\mathbf{f}
(
\mathbf{x}^*
).
}
\]

Let

\[
\mathbf{u}
=
\mathbf{x}
-
\mathbf{x}^*.
\]

Then locally,

\[
\boxed{
\mathbf{u}'
\approx
J\mathbf{u}.
}
\]

The eigenvalues of \(J\) often determine local stability.

%%%%%%%%%%%%%%%%%%%%%%%%%%%%%%%%%%%%%%%%%%%%%%%%%%%%%%%%%%%%
\subsection{Hyperbolic Equilibria}
\label{subsec:hyperbolic-equilibria-preview}
%%%%%%%%%%%%%%%%%%%%%%%%%%%%%%%%%%%%%%%%%%%%%%%%%%%%%%%%%%%%

An equilibrium is hyperbolic if the Jacobian has no eigenvalues with zero
real part.

For hyperbolic equilibria, the linearized system accurately captures the
local qualitative stability type.

If eigenvalues lie on the imaginary axis, linearization can be
inconclusive.

%%%%%%%%%%%%%%%%%%%%%%%%%%%%%%%%%%%%%%%%%%%%%%%%%%%%%%%%%%%%
\subsection{Euler's Method}
\label{subsec:eulers-method}
%%%%%%%%%%%%%%%%%%%%%%%%%%%%%%%%%%%%%%%%%%%%%%%%%%%%%%%%%%%%

For

\[
y'
=
f(t,y),
\qquad
y(t_0)=y_0,
\]

Euler's method uses step size \(h\):

\[
\boxed{
t_{n+1}
=
t_n+h,
}
\]

\[
\boxed{
y_{n+1}
=
y_n
+
h
f(t_n,y_n).
}
\]

It follows the tangent-line direction over each small step.

%%%%%%%%%%%%%%%%%%%%%%%%%%%%%%%%%%%%%%%%%%%%%%%%%%%%%%%%%%%%
\subsection{Worked Example: Euler's Method}
\label{subsec:worked-euler-method}
%%%%%%%%%%%%%%%%%%%%%%%%%%%%%%%%%%%%%%%%%%%%%%%%%%%%%%%%%%%%

\begin{workedexamplebox}{One Euler Step}

Consider

\[
y'=y,
\qquad
y(0)=1.
\]

Using step size

\[
h=0.1,
\]

Euler's method gives

\[
y_1
=
y_0
+
0.1y_0
=
1+0.1.
\]

Thus,

\begin{answerbox}
\[
\boxed{
y(0.1)
\approx
1.1.
}
\]
\end{answerbox}

The exact value is

\[
e^{0.1}
\approx
1.10517.
\]

\end{workedexamplebox}

%%%%%%%%%%%%%%%%%%%%%%%%%%%%%%%%%%%%%%%%%%%%%%%%%%%%%%%%%%%%
\subsection{Local Truncation Error}
\label{subsec:euler-local-error}
%%%%%%%%%%%%%%%%%%%%%%%%%%%%%%%%%%%%%%%%%%%%%%%%%%%%%%%%%%%%

Euler's method is based on

\[
y(t+h)
=
y(t)
+
hy'(t)
+
O(h^2).
\]

Therefore the one-step local truncation error is

\[
\boxed{
O(h^2).
}
\]

The accumulated global error over a fixed interval is typically

\[
\boxed{
O(h).
}
\]

%%%%%%%%%%%%%%%%%%%%%%%%%%%%%%%%%%%%%%%%%%%%%%%%%%%%%%%%%%%%
\subsection{Improved Euler Method Preview}
\label{subsec:improved-euler-preview}
%%%%%%%%%%%%%%%%%%%%%%%%%%%%%%%%%%%%%%%%%%%%%%%%%%%%%%%%%%%%

Improved Euler methods use more slope information.

A predictor may be

\[
\widetilde{y}_{n+1}
=
y_n
+
hf(t_n,y_n).
\]

Then a corrected value can use an average of beginning and predicted end
slopes.

This improves accuracy over ordinary Euler's method.

%%%%%%%%%%%%%%%%%%%%%%%%%%%%%%%%%%%%%%%%%%%%%%%%%%%%%%%%%%%%
\subsection{Runge--Kutta Methods Preview}
\label{subsec:runge-kutta-preview}
%%%%%%%%%%%%%%%%%%%%%%%%%%%%%%%%%%%%%%%%%%%%%%%%%%%%%%%%%%%%

Runge--Kutta methods approximate higher-order Taylor accuracy without
explicitly differentiating \(f\) repeatedly.

The classical fourth-order Runge--Kutta method uses four slope
evaluations per step and has global error of order

\[
\boxed{
O(h^4).
}
\]

These methods are developed more deeply in Chapter 11.

%%%%%%%%%%%%%%%%%%%%%%%%%%%%%%%%%%%%%%%%%%%%%%%%%%%%%%%%%%%%
\subsection{Stiff Equations Preview}
\label{subsec:stiff-equations-preview}
%%%%%%%%%%%%%%%%%%%%%%%%%%%%%%%%%%%%%%%%%%%%%%%%%%%%%%%%%%%%

Some differential equations contain dynamics occurring on very different
time scales.

Explicit numerical methods may require extremely small step sizes for
stability even when the exact solution changes slowly.

Such problems are called stiff.

Implicit numerical methods are often preferred.

%%%%%%%%%%%%%%%%%%%%%%%%%%%%%%%%%%%%%%%%%%%%%%%%%%%%%%%%%%%%
\subsection{ODE Modeling Workflow}
\label{subsec:ode-modeling-workflow}
%%%%%%%%%%%%%%%%%%%%%%%%%%%%%%%%%%%%%%%%%%%%%%%%%%%%%%%%%%%%

A practical ODE workflow is:

\begin{enumerate}
    \item define the changing quantity;
    \item identify the independent variable;
    \item express the rate law;
    \item determine initial or boundary conditions;
    \item classify the differential equation;
    \item select an analytical solution method when available;
    \item verify the solution;
    \item interpret arbitrary constants using conditions;
    \item examine equilibrium solutions;
    \item analyze qualitative behavior;
    \item use numerical methods when exact solutions are unavailable;
    \item compare predictions with observed data when modeling real
          systems;
    \item interpret parameters in the original scientific context.
\end{enumerate}

%%%%%%%%%%%%%%%%%%%%%%%%%%%%%%%%%%%%%%%%%%%%%%%%%%%%%%%%%%%%
\subsection{Choosing a First-Order Solution Method}
\label{subsec:choosing-first-order-method}
%%%%%%%%%%%%%%%%%%%%%%%%%%%%%%%%%%%%%%%%%%%%%%%%%%%%%%%%%%%%

If the equation can be written

\[
y'=g(x)h(y),
\]

use:

\[
\boxed{
\text{separation of variables}.
}
\]

If it has form

\[
y'+P(x)y=Q(x),
\]

use:

\[
\boxed{
\text{integrating factor}.
}
\]

If

\[
M\,dx+N\,dy=0
\]

and

\[
M_y=N_x,
\]

use:

\[
\boxed{
\text{exact-equation methods}.
}
\]

If

\[
y'+Py=Qy^n,
\]

consider:

\[
\boxed{
\text{Bernoulli substitution}.
}
\]

%%%%%%%%%%%%%%%%%%%%%%%%%%%%%%%%%%%%%%%%%%%%%%%%%%%%%%%%%%%%
\subsection{Common Errors}
\label{subsec:ode-common-errors}
%%%%%%%%%%%%%%%%%%%%%%%%%%%%%%%%%%%%%%%%%%%%%%%%%%%%%%%%%%%%

\subsubsection{Losing Equilibrium Solutions During Separation}

Dividing by a function of \(y\) can remove solutions where that function
equals zero.

Check these values before dividing.

%%%%%%%%%%%%%%%%%%%%%%%%%%%%%%%%%%%%%%%%%%%%%%%%%%%%%%%%%%%%
\subsubsection{Forgetting the Constant of Integration}

After indefinite integration, include an arbitrary constant.

For example,

\[
\int
\frac1y\,dy
=
\ln|y|+C.
\]

%%%%%%%%%%%%%%%%%%%%%%%%%%%%%%%%%%%%%%%%%%%%%%%%%%%%%%%%%%%%
\subsubsection{Dropping the Absolute Value in Logarithmic Integrals}

The antiderivative is

\[
\boxed{
\int
\frac1y\,dy
=
\ln|y|+C,
}
\]

not simply \(\ln y\) unless positivity is already known.

%%%%%%%%%%%%%%%%%%%%%%%%%%%%%%%%%%%%%%%%%%%%%%%%%%%%%%%%%%%%
\subsubsection{Applying Initial Conditions Too Early}

It is often cleaner to derive the general solution first and then use the
initial condition.

However, definite-integral forms can also incorporate conditions
directly when handled consistently.

%%%%%%%%%%%%%%%%%%%%%%%%%%%%%%%%%%%%%%%%%%%%%%%%%%%%%%%%%%%%
\subsubsection{Using an Integrating Factor Before Standardizing the Equation}

The equation must first be written as

\[
\boxed{
y'+P(x)y=Q(x).
}
\]

If the coefficient of \(y'\) is not \(1\), divide through first.

%%%%%%%%%%%%%%%%%%%%%%%%%%%%%%%%%%%%%%%%%%%%%%%%%%%%%%%%%%%%
\subsubsection{Forgetting the Product Rule Behind Integrating Factors}

The goal is to create

\[
\boxed{
(\mu y)'.
}
\]

One should verify that

\[
\mu'
=
P\mu.
\]

%%%%%%%%%%%%%%%%%%%%%%%%%%%%%%%%%%%%%%%%%%%%%%%%%%%%%%%%%%%%
\subsubsection{Assuming Every First-Order Equation Is Separable}

Many first-order equations cannot be rearranged into

\[
g(y)\,dy
=
f(x)\,dx.
\]

Classification should occur before algebraic manipulation.

%%%%%%%%%%%%%%%%%%%%%%%%%%%%%%%%%%%%%%%%%%%%%%%%%%%%%%%%%%%%
\subsubsection{Calling a Nonlinear Equation Linear}

Terms such as

\[
y^2,
\qquad
yy',
\qquad
\sin y
\]

make an equation nonlinear.

%%%%%%%%%%%%%%%%%%%%%%%%%%%%%%%%%%%%%%%%%%%%%%%%%%%%%%%%%%%%
\subsubsection{Ignoring the Interval of Validity}

An implicit or explicit formula may fail where denominators vanish,
logarithms become undefined, or the original differential equation loses
regularity.

Solutions should be associated with intervals of validity.

%%%%%%%%%%%%%%%%%%%%%%%%%%%%%%%%%%%%%%%%%%%%%%%%%%%%%%%%%%%%
\subsubsection{Using Superposition for Nonlinear Equations}

If \(y_1\) and \(y_2\) solve a nonlinear ODE, their sum generally does not
solve it.

Superposition is a consequence of linearity.

%%%%%%%%%%%%%%%%%%%%%%%%%%%%%%%%%%%%%%%%%%%%%%%%%%%%%%%%%%%%
\subsubsection{Forgetting Multiplicity in Characteristic Roots}

A repeated root \(r\) of multiplicity \(m\) produces

\[
e^{rx},
xe^{rx},
\ldots,
x^{m-1}e^{rx}.
\]

%%%%%%%%%%%%%%%%%%%%%%%%%%%%%%%%%%%%%%%%%%%%%%%%%%%%%%%%%%%%
\subsubsection{Forgetting the \(x\)-Multiplier Under Resonance}

If an undetermined-coefficient trial duplicates a homogeneous solution,
multiply by a sufficient power of \(x\).

%%%%%%%%%%%%%%%%%%%%%%%%%%%%%%%%%%%%%%%%%%%%%%%%%%%%%%%%%%%%
\subsubsection{Using Undetermined Coefficients for Arbitrary Forcing}

The method only applies conveniently to specific forcing families.

Variation of parameters is more general.

%%%%%%%%%%%%%%%%%%%%%%%%%%%%%%%%%%%%%%%%%%%%%%%%%%%%%%%%%%%%
\subsubsection{Confusing Stability with Boundedness}

A bounded solution is not automatically a stable equilibrium.

Stability concerns behavior of nearby solutions under perturbation.

%%%%%%%%%%%%%%%%%%%%%%%%%%%%%%%%%%%%%%%%%%%%%%%%%%%%%%%%%%%%
\subsubsection{Assuming Purely Imaginary Eigenvalues Guarantee a Nonlinear Center}

For nonlinear systems, eigenvalues with zero real part make linearization
inconclusive.

Higher-order nonlinear terms may determine stability.

%%%%%%%%%%%%%%%%%%%%%%%%%%%%%%%%%%%%%%%%%%%%%%%%%%%%%%%%%%%%
\subsubsection{Treating Euler's Approximation as Exact}

Euler's method introduces truncation error.

The approximation depends on step size.

%%%%%%%%%%%%%%%%%%%%%%%%%%%%%%%%%%%%%%%%%%%%%%%%%%%%%%%%%%%%
\subsubsection{Using a Large Euler Step Without Checking Stability}

Large step sizes can cause severe numerical error or artificial
instability even when the exact solution is stable.

%%%%%%%%%%%%%%%%%%%%%%%%%%%%%%%%%%%%%%%%%%%%%%%%%%%%%%%%%%%%
\subsection{Applications}
\label{subsec:ode-applications}
%%%%%%%%%%%%%%%%%%%%%%%%%%%%%%%%%%%%%%%%%%%%%%%%%%%%%%%%%%%%

\begin{applicationbox}

\textbf{Population Modeling.}
Exponential and logistic equations describe idealized population growth
and resource-limited growth.

\medskip

\textbf{Epidemiology.}
Compartment models such as SIR and SEIR systems use ODEs to describe
movement among disease states.

\medskip

\textbf{Physics.}
Newton's laws produce differential equations for motion, oscillation, and
forces.

\medskip

\textbf{Engineering.}
Electrical circuits, mechanical vibration, control systems, and thermal
systems are naturally represented by ODEs.

\medskip

\textbf{Chemistry.}
Reaction-rate equations describe changes in chemical concentrations over
time.

\medskip

\textbf{Environmental Science.}
ODEs model pollutant decay, ecological populations, dissolved substances,
and transport between well-mixed compartments.

\medskip

\textbf{Economics.}
Continuous-time growth, adjustment, and dynamic equilibrium models use
differential equations.

\medskip

\textbf{Operations and Applied Mathematics.}
ODEs appear in inventory models, fluid systems, queues under fluid
approximations, and many continuous-time dynamical systems.

\end{applicationbox}

%%%%%%%%%%%%%%%%%%%%%%%%%%%%%%%%%%%%%%%%%%%%%%%%%%%%%%%%%%%%
\subsection{Exercises}
\label{subsec:ode-exercises}
%%%%%%%%%%%%%%%%%%%%%%%%%%%%%%%%%%%%%%%%%%%%%%%%%%%%%%%%%%%%

\begin{exercise}[1]
Define an ordinary differential equation.
\end{exercise}

\begin{exercise}[1]
Determine the order of

\[
y'''+y'-2y=0.
\]
\end{exercise}

\begin{exercise}[1]
Determine whether

\[
y''+xy'+y=e^x
\]

is linear or nonlinear.
\end{exercise}

\begin{exercise}[1]
Determine whether

\[
y'+y^2=0
\]

is linear or nonlinear.
\end{exercise}

\begin{exercise}[1]
Define an autonomous differential equation.
\end{exercise}

\begin{exercise}[1]
Define an initial-value problem.
\end{exercise}

\begin{exercise}[1]
Define a boundary-value problem.
\end{exercise}

\begin{exercise}[1]
Define an equilibrium solution.
\end{exercise}

\begin{exercise}[1]
Define a separable equation.
\end{exercise}

\begin{exercise}[1]
State the standard form of a first-order linear ODE.
\end{exercise}

\begin{exercise}[1]
Define an exact differential equation.
\end{exercise}

\begin{exercise}[1]
Define the Wronskian.
\end{exercise}

\begin{exercise}[2]
Verify that

\[
y=Ce^{3x}
\]

solves

\[
y'=3y.
\]
\end{exercise}

\begin{exercise}[2]
Solve

\[
y'=4x.
\]
\end{exercise}

\begin{exercise}[2]
Solve

\[
y'=xy.
\]
\end{exercise}

\begin{exercise}[2]
Solve

\[
y'=-2y,
\qquad
y(0)=5.
\]
\end{exercise}

\begin{exercise}[2]
Solve

\[
y'+y=3.
\]
\end{exercise}

\begin{exercise}[2]
Find the integrating factor for

\[
y'+4xy=x.
\]
\end{exercise}

\begin{exercise}[2]
Determine whether

\[
(3x^2+2y)\,dx
+
(2x+4y^3)\,dy
=
0
\]

is exact.
\end{exercise}

\begin{exercise}[2]
Find the equilibria of

\[
y'=y(4-y).
\]
\end{exercise}

\begin{exercise}[2]
Classify the equilibria of

\[
y'=y(4-y)
\]

using a phase line.
\end{exercise}

\begin{exercise}[2]
Solve

\[
y''-3y'+2y=0.
\]
\end{exercise}

\begin{exercise}[2]
Solve

\[
y''-6y'+9y=0.
\]
\end{exercise}

\begin{exercise}[2]
Solve

\[
y''+9y=0.
\]
\end{exercise}

\begin{exercise}[2]
Find the natural frequency of

\[
4x''+100x=0.
\]
\end{exercise}

\begin{exercise}[2]
Use one Euler step with \(h=0.2\) for

\[
y'=2y,
\qquad
y(0)=1.
\]
\end{exercise}

\begin{exercise}[3]
Solve

\[
\frac{dy}{dx}
=
x(1+y^2)
\]

by separation.
\end{exercise}

\begin{exercise}[3]
Solve the IVP

\[
y'
=
2xy,
\qquad
y(0)=4.
\]
\end{exercise}

\begin{exercise}[3]
Solve

\[
y'
-
\frac2x y
=
x^2,
\qquad
x>0,
\]

using an integrating factor.
\end{exercise}

\begin{exercise}[3]
Solve

\[
(2xy+1)\,dx
+
(x^2+2y)\,dy
=
0.
\]
\end{exercise}

\begin{exercise}[3]
Transform a Bernoulli equation into a linear differential equation.
\end{exercise}

\begin{exercise}[3]
Solve a homogeneous first-order equation of the form

\[
y'
=
1+\frac{y}{x}.
\]
\end{exercise}

\begin{exercise}[3]
Derive the exponential growth solution

\[
P=P_0e^{kt}.
\]
\end{exercise}

\begin{exercise}[3]
Derive the doubling-time formula.
\end{exercise}

\begin{exercise}[3]
Derive the half-life formula.
\end{exercise}

\begin{exercise}[3]
Derive the logistic solution

\[
P(t)
=
\frac{
K
}{
1+Ae^{-rt}
}.
\]
\end{exercise}

\begin{exercise}[3]
Analyze stability of the logistic equilibria.
\end{exercise}

\begin{exercise}[3]
Derive Newton's law of cooling solution.
\end{exercise}

\begin{exercise}[3]
Formulate a mixing-tank differential equation using

\[
\text{rate in}
-
\text{rate out}.
\]
\end{exercise}

\begin{exercise}[3]
Prove the superposition principle for homogeneous linear ODEs.
\end{exercise}

\begin{exercise}[3]
Compute the Wronskian of

\[
e^x
\]

and

\[
e^{-x}.
\]
\end{exercise}

\begin{exercise}[3]
Solve

\[
y''-2y'-3y=0.
\]
\end{exercise}

\begin{exercise}[3]
Solve

\[
y''+4y'+13y=0.
\]
\end{exercise}

\begin{exercise}[3]
Solve

\[
y'''-3y''+3y'-y=0.
\]
\end{exercise}

\begin{exercise}[3]
Solve

\[
y''+y=2\cos x
\]

using undetermined coefficients, taking resonance into account.
\end{exercise}

\begin{exercise}[3]
Find a particular solution of

\[
y''-y=e^{2x}.
\]
\end{exercise}

\begin{exercise}[3]
Derive the variation-of-parameters formulas.
\end{exercise}

\begin{exercise}[3]
Use reduction of order to construct a second solution given one known
solution.
\end{exercise}

\begin{exercise}[3]
Classify a spring--mass system as over-, critically, or underdamped from
the discriminant.
\end{exercise}

\begin{exercise}[3]
Rewrite

\[
y''+3y'+2y=0
\]

as a first-order system.
\end{exercise}

\begin{exercise}[3]
Solve a \(2\times2\) diagonal linear system using eigenvalues.
\end{exercise}

\begin{exercise}[3]
Classify the origin for a planar system whose eigenvalues are

\[
-1,
\quad
-4.
\]
\end{exercise}

\begin{exercise}[3]
Classify the origin for eigenvalues

\[
2,
\quad
-3.
\]
\end{exercise}

\begin{exercise}[3]
Classify the origin for eigenvalues

\[
-1\pm3i.
\]
\end{exercise}

\begin{exercise}[3]
Derive the Euler-method update formula from a tangent-line approximation.
\end{exercise}

\begin{exercise}[3]
Explain the difference between local and global truncation error.
\end{exercise}

\begin{exercise}[4]
Prove a local existence and uniqueness theorem using Picard iteration.
\end{exercise}

\begin{exercise}[4]
Construct an IVP with more than one solution through the same initial
point.
\end{exercise}

\begin{exercise}[4]
Study maximal intervals of existence for nonlinear first-order equations.
\end{exercise}

\begin{exercise}[4]
Derive Abel's identity for the Wronskian.
\end{exercise}

\begin{exercise}[4]
Study Green's functions for linear boundary-value problems.
\end{exercise}

\begin{exercise}[4]
Develop variation of parameters for \(n\)-th order equations.
\end{exercise}

\begin{exercise}[4]
Study regular singular points and Frobenius-series solutions.
\end{exercise}

\begin{exercise}[4]
Derive the matrix exponential solution of

\[
\mathbf{x}'=A\mathbf{x}.
\]
\end{exercise}

\begin{exercise}[4]
Study Jordan-form solutions for defective matrices.
\end{exercise}

\begin{exercise}[4]
Derive the variation-of-constants formula for linear systems.
\end{exercise}

\begin{exercise}[4]
Study fundamental matrices and transition matrices.
\end{exercise}

\begin{exercise}[4]
Derive the trace--determinant classification of planar linear systems.
\end{exercise}

\begin{exercise}[4]
Study linearization near hyperbolic equilibria.
\end{exercise}

\begin{exercise}[4]
Investigate nonlinear centers and limitations of linearization.
\end{exercise}

\begin{exercise}[4]
Study Lyapunov stability methods for ODE systems.
\end{exercise}

\begin{exercise}[4]
Analyze resonance in a forced damped oscillator.
\end{exercise}

\begin{exercise}[4]
Study coupled oscillators as linear ODE systems.
\end{exercise}

\begin{exercise}[4]
Derive the local and global error orders of Euler's method.
\end{exercise}

\begin{exercise}[4]
Derive the classical fourth-order Runge--Kutta method.
\end{exercise}

\begin{exercise}[4]
Study absolute stability regions of numerical ODE methods.
\end{exercise}

\begin{exercise}[4]
Investigate stiff systems and implicit numerical methods.
\end{exercise}

%%%%%%%%%%%%%%%%%%%%%%%%%%%%%%%%%%%%%%%%%%%%%%%%%%%%%%%%%%%%
\subsection{Conceptual Summary}
\label{subsec:ode-summary}
%%%%%%%%%%%%%%%%%%%%%%%%%%%%%%%%%%%%%%%%%%%%%%%%%%%%%%%%%%%%

An ordinary differential equation relates an unknown function to its
derivatives:

\[
\boxed{
F
(
x,
y,
y',
\ldots,
y^{(n)}
)
=
0.
}
\]

A first-order autonomous equation has form

\[
\boxed{
y'=f(y).
}
\]

Equilibria satisfy

\[
\boxed{
f(y^*)=0.
}
\]

A separable equation has form

\[
\boxed{
y'=g(x)h(y).
}
\]

A first-order linear equation has form

\[
\boxed{
y'+P(x)y=Q(x),
}
\]

with integrating factor

\[
\boxed{
\mu(x)
=
e^{\int P(x)\,dx}.
}
\]

A second-order constant-coefficient homogeneous equation

\[
\boxed{
ay''+by'+cy=0
}
\]

is solved through the characteristic equation

\[
\boxed{
ar^2+br+c=0.
}
\]

For a nonhomogeneous linear equation,

\[
\boxed{
y=y_h+y_p.
}
\]

A linear system is written

\[
\boxed{
\mathbf{x}'=A\mathbf{x},
}
\]

with solution

\[
\boxed{
\mathbf{x}(t)=e^{At}\mathbf{x}_0.
}
\]

Euler's numerical method uses

\[
\boxed{
y_{n+1}
=
y_n
+
hf(t_n,y_n).
}
\]

Ordinary differential equations therefore combine

\[
\boxed{
\text{calculus}
+
\text{linear algebra}
+
\text{modeling}
+
\text{dynamical systems}
+
\text{numerical methods}.
}
\]

%%%%%%%%%%%%%%%%%%%%%%%%%%%%%%%%%%%%%%%%%%%%%%%%%%%%%%%%%%%%
\subsection{Section Connections}
\label{subsec:ode-connections}
%%%%%%%%%%%%%%%%%%%%%%%%%%%%%%%%%%%%%%%%%%%%%%%%%%%%%%%%%%%%

\chapterconnection{
Ordinary Differential Equations begins the study of continuous-time
dynamical models. First-order equations describe growth, decay, mixing,
and one-dimensional dynamics, while higher-order equations model
oscillation and mechanical systems. Linear systems connect differential
equations with eigenvalues, eigenvectors, and matrix exponentials.
Nonlinear systems lead naturally into dynamical-systems analysis, while
Euler and Runge--Kutta methods connect ODEs with numerical mathematics.
The next section, Transform Methods, develops Laplace and related
transform techniques for solving differential and integral equations,
especially problems involving initial conditions, discontinuous forcing,
and linear systems.
}

%%%%%%%%%%%%%%%%%%%%%%%%%%%%%%%%%%%%%%%%%%%%%%%%%%%%%%%%%%%%
% SECTION SOURCE TRACKING
%%%%%%%%%%%%%%%%%%%%%%%%%%%%%%%%%%%%%%%%%%%%%%%%%%%%%%%%%%%%

%------------------------------------------------------------
% 9.1 Ordinary Differential Equations
%------------------------------------------------------------
%
% Source basis:
%   - Standard ordinary differential equation theory.
%   - Standard first-order and linear ODE methods.
%   - Standard introductory dynamical-systems and numerical ODE ideas.
%
% Used for:
%   - ODE definitions and classification;
%   - order and degree;
%   - linear and nonlinear equations;
%   - autonomous and nonautonomous equations;
%   - general, particular, and singular solutions;
%   - initial-value and boundary-value problems;
%   - verification of solutions;
%   - existence and uniqueness;
%   - Picard--Lindelof theorem;
%   - direction fields;
%   - equilibrium and phase-line analysis;
%   - separable equations;
%   - first-order linear equations;
%   - integrating factors;
%   - exact equations;
%   - Bernoulli equations;
%   - homogeneous substitutions;
%   - exponential growth and decay;
%   - logistic growth;
%   - Newton cooling;
%   - mixing problems;
%   - second-order linear equations;
%   - superposition;
%   - Wronskians;
%   - Abel identity preview;
%   - characteristic equations;
%   - distinct, repeated, and complex roots;
%   - higher-order constant-coefficient equations;
%   - nonhomogeneous equations;
%   - undetermined coefficients;
%   - resonance;
%   - variation of parameters;
%   - reduction of order;
%   - mechanical vibrations;
%   - damping and forced oscillation;
%   - systems of first-order ODEs;
%   - eigenvalue methods;
%   - matrix exponentials;
%   - fundamental matrices;
%   - planar linear systems;
%   - trace-determinant classification;
%   - nonlinear systems;
%   - linearization;
%   - Euler's method;
%   - improved Euler and Runge--Kutta previews;
%   - stiffness preview;
%   - applications and exercises.
%
% Notes:
%   - Transform Methods is developed next in Section 9.2.
%   - Dynamical Systems is developed later in Chapter 9.
%   - Numerical ODE methods are developed more deeply in Chapter 11.
%
%%%%%%%%%%%%%%%%%%%%%%%%%%%%%%%%%%%%%%%%%%%%%%%%%%%%%%%%%%%%